\documentclass[11pt]{article}
% ------------------------------------------------------------------
% Preámbulo común a la solución del taller y a la guía de estudio
% ------------------------------------------------------------------
\usepackage[T1]{fontenc}
\usepackage[utf8]{inputenc}
\usepackage[spanish,es-nodecimaldot,es-noshorthands,es-tabla,es-lcroman]{babel}
\usepackage{mathpazo}
\usepackage{tgpagella}
\usepackage[scaled=.92]{helvet}
\renewcommand{\ttdefault}{lmtt}
\linespread{1.06}
\usepackage{amsmath,amssymb,amsthm,mathtools}
\usepackage[letterpaper,margin=2.25cm,headheight=14pt]{geometry}
\usepackage[table,dvipsnames]{xcolor}
\usepackage{array,booktabs,longtable,multicol,graphicx}
\usepackage{enumitem}
\usepackage{microtype}
\usepackage{fancyhdr}
\usepackage{titlesec}
\usepackage{tikz}
\usetikzlibrary{arrows.meta,positioning,calc,shapes.geometric,fit,backgrounds}
\usepackage[most]{tcolorbox}
\usepackage[hidelinks,unicode]{hyperref}
\usepackage{bookmark}

% ----------------------------------------------------------- colores
\definecolor{tinta}{RGB}{24,52,94}
\definecolor{tinta2}{RGB}{46,96,150}
\definecolor{verde}{RGB}{22,102,82}
\definecolor{ambar}{RGB}{160,96,8}
\definecolor{rojo}{RGB}{160,36,36}
\definecolor{malo}{RGB}{250,214,214}
\definecolor{bueno}{RGB}{214,236,224}
\definecolor{bueno2}{RGB}{218,228,244}
\definecolor{papel}{RGB}{246,248,251}
\hypersetup{colorlinks=true,linkcolor=tinta2,urlcolor=tinta2,citecolor=tinta2}

% ---------------------------------------------------------- formato
\setlength{\parindent}{0pt}
\setlength{\parskip}{0.55em}
\setlist{itemsep=0.2em,topsep=0.3em,parsep=0.15em}
\titleformat{\section}{\Large\bfseries\sffamily\color{tinta}}{\thesection}{0.6em}{}[\vspace{-0.4em}{\color{tinta!35}\rule{\linewidth}{0.8pt}}]
\titleformat{\subsection}{\large\bfseries\sffamily\color{tinta}}{\thesubsection}{0.6em}{}
\titleformat{\subsubsection}{\normalsize\bfseries\sffamily\color{tinta2}}{\thesubsubsection}{0.6em}{}
\titlespacing*{\section}{0pt}{1.4em}{0.6em}
\titlespacing*{\subsection}{0pt}{1.0em}{0.3em}
\titlespacing*{\subsubsection}{0pt}{0.8em}{0.2em}

% ----------------------------------------------------------- macros
\newcommand{\s}[1]{\{#1\}}                         % conjunto por extensión
\newcommand{\vac}{\varnothing}
\newcommand{\cl}[1]{\overline{#1}}                 % clausura
\newcommand{\inte}[1]{\mathring{#1}}               % interior (letra sola)
\newcommand{\Int}{\operatorname{int}}
\newcommand{\Ext}{\operatorname{ext}}
\newcommand{\Fr}{\operatorname{Fr}}
\newcommand{\cotau}{\mathrm{co}\text{-}\tau}
\newcommand{\co}[1]{\mathrm{co}\text{-}#1}
\newcommand{\N}{\mathcal{N}}
\newcommand{\B}{\mathcal{B}}
\newcommand{\C}{\mathfrak{C}}
\newcommand{\Sb}{\mathcal{S}}
\newcommand{\Pt}{\mathcal{P}}
\newcommand{\R}{\mathbb{R}}
\newcommand{\Z}{\mathbb{Z}}
\newcommand{\Q}{\mathbb{Q}}
\newcommand{\sm}{\smallsetminus}
\newcommand{\homeo}{\cong}
\newcommand{\si}{\textcolor{verde}{\textbf{sí}}}
\newcommand{\no}{\textcolor{rojo}{\textbf{no}}}
\newcommand{\ok}{\textcolor{verde}{$\checkmark$}}
\newcommand{\mal}{\textcolor{rojo}{$\times$}}
\newcommand{\fuente}[1]{\textsuperscript{\textcolor{tinta2}{[#1]}}}

% \input expandible, para poder usarlo dentro de tablas antes de \bottomrule
\makeatletter\newcommand{\tinput}[1]{\@@input gen/#1.tex }\makeatother

% ------------------------------------------------------------ cajas
\tcbset{base/.style={enhanced,breakable,boxrule=0pt,arc=1.5pt,left=9pt,right=9pt,top=6pt,bottom=6pt,
  before skip=0.8em,after skip=0.8em,fonttitle=\bfseries\sffamily\small,coltitle=white,
  attach boxed title to top left={xshift=6pt,yshift=-2.5mm},boxed title style={arc=1.5pt,boxrule=0pt,top=1.5pt,bottom=1.5pt}}}
\newtcolorbox{definicion}[1][Definición]{base,colback=tinta!5,borderline west={2.5pt}{0pt}{tinta},
  colbacktitle=tinta,title={#1}}
\newtcolorbox{teorema}[1][Teorema]{base,colback=tinta2!6,borderline west={2.5pt}{0pt}{tinta2},
  colbacktitle=tinta2,title={#1}}
\newtcolorbox{ejemplo}[1][Ejemplo]{base,colback=verde!5,borderline west={2.5pt}{0pt}{verde},
  colbacktitle=verde,title={#1}}
\newtcolorbox{ojo}[1][Advertencia]{base,colback=ambar!7,borderline west={2.5pt}{0pt}{ambar},
  colbacktitle=ambar,title={#1}}
\newtcolorbox{metodo}[1][Método]{base,colback=papel,borderline west={2.5pt}{0pt}{black!55},
  colbacktitle=black!65,title={#1}}
\newtcolorbox{respuesta}{enhanced,boxrule=0.6pt,colframe=verde!70!black,colback=verde!4,arc=1.5pt,
  left=8pt,right=8pt,top=4pt,bottom=4pt,before skip=0.6em,after skip=0.8em}
\newcommand{\resp}[1]{\begin{respuesta}\textbf{\sffamily\color{verde!70!black}Respuesta.} #1\end{respuesta}}
\newtcolorbox{enunciado}{enhanced,breakable,boxrule=0pt,colback=black!4,arc=1.5pt,left=8pt,right=8pt,top=5pt,bottom=5pt,
  before skip=0.6em,after skip=0.6em,fontupper=\itshape}

% -------------------------------------------------- estilos de TikZ
\tikzset{
  hasse/.style={draw=tinta!45,line width=0.35pt},
  nodoh/.style={circle,draw=tinta,fill=white,inner sep=0.8pt,minimum size=5.6mm,font=\scriptsize},
  pto/.style={inner sep=2pt,font=\normalsize},
  flecha/.style={-{Stealth[length=2.2mm]},line width=0.55pt},
  ovalo/.style={draw=black!70,line width=0.5pt,ellipse,minimum width=2.1cm,minimum height=3.9cm},
}
% diagrama de flechas de una función entre conjuntos de 4 puntos:
% \diagrama{nombre}{d1,d2,d3,d4}{c1,c2,c3,c4}{i1/j1,i2/j2,...}
\newcommand{\diagrama}[4]{%
\begin{tikzpicture}[baseline=(current bounding box.center)]
  \node[ovalo] (L) at (0,0) {};
  \node[ovalo] (Rr) at (3.6,0) {};
  \node at (1.8,2.05) {$#1$};
  \foreach \t [count=\i] in {#2} {\node[pto] (l\i) at (0,1.5-\i*0.6) {$\t$};}
  \foreach \t [count=\i] in {#3} {\node[pto] (r\i) at (3.6,1.5-\i*0.6) {$\t$};}
  \foreach \a/\b in {#4} {\draw[flecha] (l\a) -- (r\b);}
\end{tikzpicture}}

\hypersetup{pdftitle={Guía de estudio — Taller 2 de Topología},pdfsubject={Topología general},pdfkeywords={topología, guía de estudio, parcial}}
\pagestyle{fancy}
\fancyhf{}
\fancyhead[L]{\small\sffamily\color{black!60}Guía de estudio — Topología, módulo 2}
\fancyhead[R]{\small\sffamily\color{black!60}\nouppercase{\leftmark}}
\fancyfoot[C]{\small\sffamily\thepage}
\renewcommand{\headrulewidth}{0.3pt}
\renewcommand{\sectionmark}[1]{\markboth{#1}{}}
\newcommand{\paresTotal}{406}\newcommand{\paresComp}{163}\newcommand{\paresNoComp}{243}


\begin{document}
\thispagestyle{empty}
\vspace*{1.2cm}
\begin{center}
{\sffamily\large UNIVERSIDAD DE PANAMÁ · FACINET · LICENCIATURA EN MATEMÁTICA}\\[2.2em]
{\sffamily\bfseries\Huge\color{tinta} Guía de estudio}\\[0.5em]
{\sffamily\bfseries\LARGE\color{tinta} Topología — Módulo 2}\\[1.0em]
{\sffamily\Large Todo lo que hay detrás del Taller 2, explicado paso a paso para el parcial}\\[2.2em]
\begin{tikzpicture}[x=1cm,y=1cm,
  caja/.style={draw=tinta,fill=tinta!6,rounded corners=2pt,font=\sffamily\small,align=center,inner sep=5pt,minimum width=3.3cm,minimum height=1.05cm},
  fl/.style={-{Stealth[length=2.2mm]},draw=tinta!70,line width=0.6pt}]
\node[caja] (top) at (0,0) {Topología $\tau$\\{\scriptsize abiertos}};
\node[caja] (cer) at (-4.6,-1.9) {Cerrados\\{\scriptsize $\cotau$}};
\node[caja] (vec) at (0,-1.9) {Vecindades\\{\scriptsize $\N_x$}};
\node[caja] (bas) at (4.6,-1.9) {Bases y sub-bases\\{\scriptsize generar $\tau$}};
\node[caja] (der) at (-4.6,-3.9) {Derivado y clausura\\{\scriptsize $A'$, $\cl{A}$}};
\node[caja] (int) at (0,-3.9) {Interior, exterior,\\frontera};
\node[caja] (sub) at (4.6,-3.9) {Subespacios y\\comparación};
\node[caja] (con) at (-2.3,-5.9) {Continuidad;\\funciones abiertas y cerradas};
\node[caja] (hom) at (3.2,-5.9) {Homeomorfismos\\{\scriptsize $X\homeo Y$}};
\draw[fl] (top) -- (cer); \draw[fl] (top) -- (vec); \draw[fl] (top) -- (bas);
\draw[fl] (cer) -- (der); \draw[fl] (vec) -- (der); \draw[fl] (vec) -- (int); \draw[fl] (bas) -- (sub);
\draw[fl] (der) -- (con); \draw[fl] (int) -- (con); \draw[fl] (sub) -- (hom); \draw[fl] (con) -- (hom);
\end{tikzpicture}
\end{center}
\vspace{1.2em}
\begin{metodo}[Cómo usar esta guía]
\begin{itemize}
\item Cada sección sigue el mismo orden: \textbf{idea} en lenguaje llano, \textbf{definición} precisa, \textbf{resultados}
con su demostración cuando es corta e instructiva, \textbf{método} de cálculo paso a paso y \textbf{ejemplos}, casi todos
tomados del Taller~2 para que teoría y práctica coincidan.
\item Las cajas ámbar marcan los errores más comunes y los puntos donde las convenciones cambian de un libro a otro.
\item La sección~\ref{sec:recetario} es un recetario de una página para repasar la víspera; la~\ref{sec:auto} es una
autoevaluación con respuestas.
\item Los números entre corchetes remiten a las fuentes de la sección final. Todos los cálculos con conjuntos finitos
de esta guía se comprobaron por computadora con el mismo programa que verificó el taller.
\end{itemize}
\end{metodo}
\newpage
\setcounter{tocdepth}{2}
{\small\tableofcontents}
\newpage

\section{Espacios topológicos}

\subsection{La idea: hablar de «cercanía» sin medir distancias}
En cálculo, casi todo (límites, continuidad, convergencia) se formula con desigualdades del tipo $|x-a|<\varepsilon$:
«$x$ está cerca de $a$». La observación que da origen a la topología es que no hace falta la distancia en sí; basta
saber cuáles conjuntos son \emph{abiertos}. En $\R$, un conjunto $U$ es abierto si cada uno de sus puntos tiene un
pequeño intervalo alrededor contenido en $U$; es decir, si cada punto de $U$ está «rodeado» por puntos de $U$.
Con eso solo ya se puede decir qué es una función continua (sección~\ref{sec:cont}).

La topología general da el paso siguiente: en un conjunto cualquiera $X$, \emph{declarar} cuáles subconjuntos se
llamarán abiertos, exigiendo únicamente tres reglas que los abiertos de $\R$ cumplen. Los ejemplos finitos del taller
son el laboratorio ideal: son tan pequeños que todo se puede calcular a mano, y obligan a usar las definiciones en
lugar de la intuición geométrica.

\subsection{Definición}
\begin{definicion}[Definición: topología y espacio topológico]
Sea $X$ un conjunto. Una \textbf{topología} sobre $X$ es una familia $\tau$ de subconjuntos de $X$ tal que:
\begin{description}[leftmargin=1.4cm,labelwidth=1.0cm,itemsep=0.1em]
\item[(T1)] $\vac\in\tau$ y $X\in\tau$;
\item[(T2)] la unión de \emph{cualquier} subfamilia de $\tau$ (finita o infinita) pertenece a $\tau$;
\item[(T3)] la intersección de \emph{dos} miembros de $\tau$ (y por inducción, de un número finito) pertenece a $\tau$.
\end{description}
El par $(X,\tau)$ se llama \textbf{espacio topológico} y los miembros de $\tau$ son sus \textbf{conjuntos abiertos}.
\end{definicion}
Es la definición de todos los textos de referencia\fuente{1,2,3}; algunos omiten (T1) porque consideran que la unión de
la familia vacía es $\vac$ y su intersección es $X$.

\begin{ojo}[¿Por qué uniones arbitrarias pero solo intersecciones finitas?]
Porque así ocurre en $\R$. Los intervalos $U_n=(-\tfrac1n,\tfrac1n)$ son abiertos, pero su intersección
$\bigcap_{n\ge1}U_n=\{0\}$ no lo es. Si se exigiera cerrar bajo intersecciones infinitas, la topología usual de $\R$
no sería una topología. En un espacio \emph{finito} la distinción desaparece: solo hay finitos abiertos, y toda
intersección es finita.
\end{ojo}

\subsection{Los ejemplos que hay que tener en la cabeza}
\begin{center}\small\renewcommand{\arraystretch}{1.3}
\begin{tabular}{@{}p{3.6cm}p{5.6cm}p{6.2cm}@{}}\toprule
\textbf{Nombre} & \textbf{Abiertos} & \textbf{Comentario}\\ \midrule
Indiscreta (trivial) & $\{\vac,X\}$ & La menor topología posible: la más gruesa.\\
Discreta & $\Pt(X)$: todos los subconjuntos & La mayor posible: la más fina. Todo punto es abierto.\\
Usual de $\R$ & uniones de intervalos abiertos $(a,b)$ & La del cálculo. Problema IV.2.\\
Del punto incluido & $\vac$ y los conjuntos que contienen a un punto fijo $p$ & Problema III.7, con $X=\R$, $p=0$\fuente{10,18}.\\
De partición & uniones de bloques de una partición de $X$ & Problema III.6, con bloques $\{2n-1,2n\}$\fuente{10,19}.\\
Cofinita & $\vac$ y los conjuntos de complemento finito & No aparece en el taller; es el contraejemplo clásico siguiente.\\
Cadena & $\vac\subset A_1\subset A_2\subset\dots\subset X$ & Cualquier cadena que incluya $\vac$ y $X$ es topología. Problemas III.1 a III.3.\\
Sierpiński & $X=\{0,1\}$, $\tau=\{\vac,\{1\},X\}$ & El espacio no trivial más pequeño.\\ \bottomrule
\end{tabular}\end{center}

\begin{ejemplo}[Ejemplo: por qué una topología de partición es topología]
Sea $\{P_i\}_{i\in I}$ una partición de $X$ y $\tau$ la familia de las uniones de bloques (incluida la unión vacía,
$\vac$). (T1) $\vac$ es la unión de ningún bloque y $X$ la de todos. (T2) Unión de uniones de bloques es unión de
bloques. (T3) Si $A=\bigcup_{i\in I_A}P_i$ y $B=\bigcup_{i\in I_B}P_i$, entonces $A\cap B=\bigcup_{i\in I_A\cap I_B}P_i$,
porque bloques distintos son disjuntos. En el problema III.6 se hizo la misma demostración, escrita con la condición
«$2n\in A\iff 2n-1\in A$». En estas topologías todo abierto es también cerrado.
\end{ejemplo}

\subsection{Cómo verificar si una familia finita es topología}\label{sec:verif}
\begin{metodo}[Método: verificación en el caso finito]
\begin{enumerate}[leftmargin=1.6em]
\item Comprobar que $\vac$ y $X$ están en la lista.
\item Apartar $\vac$ y $X$. Con los $m$ conjuntos restantes, formar la tabla de uniones y la tabla de intersecciones
($m\times m$; como son simétricas basta la mitad).
\item Si todos los resultados están en la lista, \textbf{es topología}. Si alguno falta, \textbf{no lo es}, y ese par
es el contraejemplo que hay que escribir.
\end{enumerate}
\end{metodo}
\emph{Por qué basta de dos en dos.} La familia es finita, así que toda unión de miembros es una unión de finitos
conjuntos $A_1\cup\dots\cup A_n$, que se calcula como $(\dots((A_1\cup A_2)\cup A_3)\dots)\cup A_n$: si cada paso se
queda dentro de la familia, el resultado también. Lo mismo para intersecciones.

\emph{Atajos que ahorran trabajo.} (1) Las operaciones con $\vac$ y con $X$ nunca producen nada nuevo.
(2) Si $A\subseteq B$, entonces $A\cup B=B$ y $A\cap B=A$: los pares encajados se cumplen solos. Solo hay que mirar
los pares \emph{no comparables}. (3) Si la familia tiene una descripción estructural (por ejemplo «todos los
subconjuntos de $\{1,2,6\}$, más $X$», como $\tau_1$ de la Parte~II), el argumento general es más rápido que la tabla.

\begin{ejemplo}[Ejemplo: Parte I del taller]
$X=\s{a,b,c,d,e}$.
\begin{itemize}
\item $\tau_1=\{X,\vac,\s{a},\s{c,d},\s{a,c,d},\s{b,c,d,e}\}$. Pares no comparables: $\s{a}$ con $\s{c,d}$ (unión
$\s{a,c,d}$, intersección $\vac$), $\s{a}$ con $\s{b,c,d,e}$ (unión $X$, intersección $\vac$), $\s{a,c,d}$ con
$\s{b,c,d,e}$ (unión $X$, intersección $\s{c,d}$). Todo está en $\tau_1$: \textbf{es topología}.
\item $\tau_2$ cambia $\s{b,c,d,e}$ por $\s{b,c,d}$: ahora $\s{a}\cup\s{b,c,d}=\s{a,b,c,d}\notin\tau_2$.
\textbf{Falla (T2).}
\item $\tau_3$ lo cambia por $\s{a,b,d,e}$: $\s{c,d}\cap\s{a,b,d,e}=\s{d}\notin\tau_3$. \textbf{Falla (T3).}
\end{itemize}
Moraleja para el examen: un solo par que falle basta para concluir que no es topología, pero para afirmar que
\emph{sí} lo es hay que revisar todos los pares no comparables.
\end{ejemplo}

\begin{ojo}[Errores frecuentes al verificar]
\begin{itemize}
\item Revisar solo las uniones (o solo las intersecciones). $\tau_3$ pasa todas las uniones y aun así no es topología.
\item Olvidar que $\vac$ y $X$ deben estar.
\item Confundir «el conjunto $\s{a,b}$ no está en $\tau$» con «los puntos $a$ y $b$ no están cubiertos». Lo que
pertenece a $\tau$ son subconjuntos, no puntos.
\end{itemize}
\end{ojo}

\subsection{¿Cuántas topologías hay en un conjunto finito?}\label{sec:cuantas}
No se conoce una fórmula cerrada. Los primeros valores son\fuente{11,12,13,21}:
\begin{center}\small\renewcommand{\arraystretch}{1.2}
\begin{tabular}{@{}lcccccc@{}}\toprule
número de puntos $n$ & 0 & 1 & 2 & 3 & 4 & 5\\ \midrule
topologías distintas (OEIS A000798) & 1 & 1 & 4 & 29 & 355 & 6942\\
salvo homeomorfismo (OEIS A001930) & 1 & 1 & 3 & 9 & 33 & 139\\ \bottomrule
\end{tabular}\end{center}
Los valores de las dos filas para $n\le5$ se reprodujeron por computadora (contando preórdenes, que en un conjunto
finito corresponden uno a uno con las topologías\fuente{13,21}); coinciden con los de la OEIS. Para $n=3$, las 29 topologías se reparten por número de abiertos así:
\begin{center}\small\renewcommand{\arraystretch}{1.2}
\begin{tabular}{@{}lcccccc|c@{}}\toprule
número de abiertos & 2 & 3 & 4 & 5 & 6 & 8 & total\\ \midrule
cuántas topologías & 1 & 6 & 9 & 6 & 6 & 1 & 29\\ \bottomrule
\end{tabular}\end{center}
Nótese que no hay ninguna con 7 abiertos. Para $n=4$ (355 en total): 1 con 2 abiertos, 14 con 3, 43 con 4, 60 con 5,
72 con 6, 54 con 7, 54 con 8, 20 con 9, 24 con 10, 12 con 12 y 1 con 16.

\begin{teorema}[Topologías con pocos abiertos (problemas III.1 a III.3 y III.13)]
Sea $X$ un conjunto y $A$, $B$, $C$ subconjuntos propios y no vacíos, distintos entre sí.
\begin{enumerate}[label=(\alph*),leftmargin=1.9em]
\item $\{\vac,A,X\}$ siempre es topología.
\item $\{\vac,A,B,X\}$ es topología $\iff$ $A\subset B$, o $B\subset A$, o $B=X\sm A$.
\item $\{\vac,A,B,C,X\}$ es topología $\iff$ (salvo reordenar) es una cadena $A\subset B\subset C$; o
$A\cap B=\vac$ y $C=A\cup B$; o $A\cup B=X$ y $C=A\cap B$.
\end{enumerate}
\end{teorema}
\emph{Demostración de (b).} ($\Leftarrow$) En los tres casos $A\cup B$ y $A\cap B$ son miembros. ($\Rightarrow$) Si $A$
y $B$ no son comparables, $A\cup B$ no puede ser $A$, ni $B$, ni $\vac$; solo queda $A\cup B=X$. Y $A\cap B$ no puede
ser $A$, ni $B$, ni $X$; solo queda $A\cap B=\vac$. Luego $B=X\sm A$. \qed

\emph{Idea de (c).} Si los tres son comparables dos a dos, forman una cadena. Si dos de ellos, digamos $A$ y $B$, no
son comparables, entonces $A\cup B$ y $A\cap B$ son miembros distintos de $A$ y de $B$; como solo quedan $\vac$, $C$ y
$X$ disponibles, o bien $A\cap B=\vac$ y $A\cup B=C$, o bien $A\cap B=C$ y $A\cup B=X$ (el caso $A\cap B=\vac$ y
$A\cup B=X$ no deja lugar para $C$: se comprueba que $C$ no podría ser comparable con ambos sin crear un conjunto
nuevo).

\begin{ejemplo}[Ejemplo: contar con el teorema]
En $X=\s{2,4,6}$ (problema III.1): cadenas $\s{x}\subset\s{x,y}$: $3\cdot2=6$; complementarios $\s{x}$ y
$X\sm\s{x}$: $3$. Total $9$, como dice la tabla anterior. En un conjunto de 4 puntos: $36$ cadenas y $7$ pares
complementarios, total $43$.
\end{ejemplo}

% =====================================================================
\section{Conjuntos cerrados}
\begin{definicion}[Definición: cerrado]
Un conjunto $F\subseteq X$ es \textbf{cerrado} si su complemento $X\sm F$ es abierto. La familia de los cerrados se
denota aquí $\cotau$ (en el taller, «$c_0\tau$»).
\end{definicion}
\begin{teorema}[Propiedades de los cerrados]
\begin{enumerate}[label=(C\arabic*),leftmargin=2.6em]
\item $\vac$ y $X$ son cerrados.
\item La intersección de \emph{cualquier} familia de cerrados es cerrada.
\item La unión de \emph{dos} (o de finitos) cerrados es cerrada.
\end{enumerate}
\end{teorema}
\emph{Demostración.} Son los axiomas (T1) a (T3) pasados por las leyes de De Morgan:
$X\sm\bigcap_i F_i=\bigcup_i (X\sm F_i)$ y $X\sm(F_1\cup F_2)=(X\sm F_1)\cap(X\sm F_2)$. \qed

Obsérvese el intercambio: los abiertos toleran uniones arbitrarias e intersecciones finitas; los cerrados, al revés.

\begin{metodo}[Método: hallar $\cotau$]
Tomar el complemento de cada abierto, uno por uno. La lista de cerrados tiene exactamente tantos elementos como la
de abiertos (el complemento es una biyección).
\end{metodo}

\begin{ojo}[«Cerrado» no es lo contrario de «abierto»]
Un conjunto puede ser:
\begin{itemize}
\item \textbf{abierto y cerrado a la vez}: siempre lo son $\vac$ y $X$; en $\tau_1$ de la Parte~I también $\s{a}$ y
$\s{b,c,d,e}$ (cada uno es el complemento del otro);
\item \textbf{ni abierto ni cerrado}: en $\tau_1$, el conjunto $\s{b}$ no es abierto, y tampoco es cerrado porque su
complemento $\s{a,c,d,e}$ no está en $\tau_1$; en $\R$, el intervalo $[0,1)$;
\item solo abierto, o solo cerrado.
\end{itemize}
Por eso, para probar que algo \emph{no} es cerrado no sirve decir «es abierto»: hay que mostrar que su complemento no
es abierto.
\end{ojo}

\begin{ejemplo}[Ejemplo: los cerrados en el taller]
\begin{itemize}
\item Parte I, $\tau_1$: $\co{\tau_1}=\{\vac,X,\s{a},\s{b,e},\s{a,b,e},\s{b,c,d,e}\}$.
\item Parte I, $\tau_4$: $\co{\tau_4}=\{\vac,X,\s{a},\s{c,d},\s{a,c,d},\s{b,c,d,e}\}$, que es exactamente $\tau_1$. No es
casualidad: en un espacio finito la familia de cerrados cumple también los axiomas de topología (finitas uniones e
intersecciones), así que $\cotau$ es otra topología sobre $X$.
\item Parte II, $\tau_2$: $\co{\tau_2}=\{\vac,X,\s{1},\s{1,3},\s{1,4,6},\s{1,2,4,6},\s{1,3,4,6}\}$.
\end{itemize}
\end{ejemplo}

\section{Vecindades}
\subsection{Idea y definición}
Una vecindad de $x$ es un conjunto que contiene «todo lo que está suficientemente cerca de $x$». En $\R$: un conjunto
que contiene algún intervalo $(x-\varepsilon,x+\varepsilon)$.
\begin{definicion}[Definición: vecindad y sistema de vecindades]
Sea $(X,\tau)$ un espacio y $x\in X$. Un conjunto $N\subseteq X$ es una \textbf{vecindad} (o \emph{entorno}) de $x$ si
existe un abierto $U$ tal que
\[x\in U\subseteq N.\]
La familia de todas las vecindades de $x$ se llama \textbf{sistema de vecindades} de $x$ y se denota $\N_x$.
\end{definicion}
\begin{ojo}[Dos convenciones]
En Lipschutz, Kelley, Willard y Bourbaki\fuente{1,3,4,26} la vecindad \emph{no} tiene que ser abierta: basta que
contenga un abierto que contenga al punto. En Munkres\fuente{2}, «vecindad de $x$» significa directamente «abierto que
contiene a $x$». Los enunciados del taller («encuentre $N_c$») corresponden a la primera. Si en el parcial se usa la
segunda, la respuesta es la sublista de vecindades que son abiertas. Ante la duda, conviene escribir ambas.
\end{ojo}

\begin{teorema}[Propiedades del sistema de vecindades]
\begin{enumerate}[label=(N\arabic*),leftmargin=2.6em]
\item $\N_x\neq\vac$ (porque $X\in\N_x$) y todo $N\in\N_x$ contiene a $x$.
\item Si $N\in\N_x$ y $N\subseteq M$, entonces $M\in\N_x$ (los superconjuntos de vecindades son vecindades).
\item Si $N,M\in\N_x$, entonces $N\cap M\in\N_x$.
\item Un conjunto $U$ es abierto $\iff$ $U$ es vecindad de cada uno de sus puntos.
\end{enumerate}
\end{teorema}
\emph{Demostración de (N4).} ($\Rightarrow$) Si $U$ es abierto y $x\in U$, entonces $x\in U\subseteq U$. ($\Leftarrow$)
Para cada $x\in U$ hay un abierto $U_x$ con $x\in U_x\subseteq U$; entonces $U=\bigcup_{x\in U}U_x$ es unión de
abiertos. \qed

La propiedad (N4) es la que más se usa: dice que para probar que un conjunto es abierto basta «rodear» cada uno de
sus puntos con un abierto contenido en él.

\subsection{El abierto mínimo: la herramienta clave en espacios finitos}\label{sec:minimo}
\begin{teorema}[Abierto mínimo]
Si $\tau$ es finita, para cada $x\in X$ el conjunto
\[U_x=\bigcap\{U\in\tau:\ x\in U\}\]
es abierto y es el menor abierto que contiene a $x$. En consecuencia $\N_x=\{N\subseteq X:\ U_x\subseteq N\}$, y si
$X$ tiene $n$ puntos y $U_x$ tiene $k$, entonces $\N_x$ tiene $2^{\,n-k}$ elementos.
\end{teorema}
\emph{Demostración.} $U_x$ es intersección de finitos abiertos, luego es abierto, contiene a $x$ y está contenido en
todo abierto que contiene a $x$. Si $U_x\subseteq N$, $N$ es vecindad. Si $N$ es vecindad, $x\in U\subseteq N$ para
algún abierto $U$, y $U_x\subseteq U\subseteq N$. Las vecindades se obtienen agregando a $U_x$ cualquier subconjunto de
los $n-k$ puntos restantes. \qed

En espacios infinitos el abierto mínimo puede no existir: en $\R$ usual, la intersección de todos los abiertos que
contienen a $0$ es $\{0\}$, que no es abierto. Los espacios en que siempre existe se llaman \emph{de Alexandrov}; todo
espacio finito lo es, y también las topologías de los problemas III.6 y III.7\fuente{13,14}.

\begin{metodo}[Método: calcular $U_x$ y $\N_x$]
\begin{enumerate}[leftmargin=1.6em]
\item Listar los abiertos que contienen a $x$.
\item Intersecarlos: eso es $U_x$. (Suele bastar con mirar cuál es el más pequeño de la lista; si hay dos no
comparables, su intersección también está en la lista por (T3).)
\item $\N_x$ consta de $U_x$ y de todos los conjuntos que resultan de agregarle puntos.
\end{enumerate}
\end{metodo}

\begin{ejemplo}[Ejemplo: Parte I, $\tau_1=\{X,\vac,\s{a},\s{c,d},\s{a,c,d},\s{b,c,d,e}\}$]
\begin{center}\small\renewcommand{\arraystretch}{1.15}
\begin{tabular}{@{}c|ccccc@{}}
$x$ & $a$ & $b$ & $c$ & $d$ & $e$\\ \hline
$U_x$ & $\s{a}$ & $\s{b,c,d,e}$ & $\s{c,d}$ & $\s{c,d}$ & $\s{b,c,d,e}$
\end{tabular}\end{center}
$\N_c$: a $\s{c,d}$ se le agrega cualquier subconjunto de $\s{a,b,e}$: $2^3=8$ vecindades. $\N_e$: a $\s{b,c,d,e}$
solo se le puede agregar $a$: $2$ vecindades, $\s{b,c,d,e}$ y $X$.

Dos lecturas útiles de la tabla: (1) $c$ y $d$ tienen el mismo abierto mínimo: la topología «no los distingue»; lo
mismo $b$ y $e$. (2) $U_b$ contiene a $c$ pero $U_c$ no contiene a $b$: la relación «estar en el abierto mínimo de» no
es simétrica.
\end{ejemplo}
\begin{ejemplo}[Ejemplo: vecindades en los problemas III.6 y III.7]
En $\Z$ con la topología par--impar, $U_5=\s{5,6}$: todo abierto que contiene a $5=2\cdot3-1$ debe contener a
$6=2\cdot3$, y $\s{5,6}$ es abierto. Así $\N_5=\{N\subseteq\Z:\s{5,6}\subseteq N\}$.

En $\R$ con la topología del punto incluido $0$, $U_1=\s{0,1}$: todo abierto no vacío contiene a $0$. Así
$\N_1=\{N\subseteq\R:\s{0,1}\subseteq N\}$. En general $U_x=\s{0,x}$ y $U_0=\s{0}$.
\end{ejemplo}

% =====================================================================
\section{Puntos de acumulación y conjunto derivado}
\subsection{Definición, leída con lupa}
\begin{definicion}[Definición: punto de acumulación y derivado]
Sea $A\subseteq X$. Un punto $x\in X$ es \textbf{punto de acumulación} (o punto límite) de $A$ si \emph{todo} abierto
$U$ que contiene a $x$ contiene algún punto de $A$ \emph{distinto de $x$}:
\[\forall\,U\in\tau\ \big(x\in U\ \Longrightarrow\ (U\sm\{x\})\cap A\neq\vac\big).\]
El conjunto de los puntos de acumulación de $A$ es el \textbf{conjunto derivado} de $A$, denotado $A'$.
\end{definicion}
Tres detalles de la definición que deciden los ejercicios:
\begin{enumerate}[leftmargin=1.6em]
\item \textbf{«Todo abierto».} Para probar que $x\in A'$ hay que revisar todos los abiertos que contienen a $x$; para
probar que $x\notin A'$ basta exhibir \emph{uno} que falle.
\item \textbf{«Distinto de $x$».} El propio $x$ no cuenta. Por eso un punto de $A$ puede no ser de acumulación de $A$
(se llama \textbf{punto aislado} de $A$), y un punto que no está en $A$ sí puede serlo.
\item \textbf{$x$ recorre todo $X$}, no solo $A$. Hay que examinar cada punto del espacio.
\end{enumerate}

\begin{ejemplo}[Ejemplos en $\R$ con la topología usual]
$(0,1)'=[0,1]$: los extremos $0$ y $1$ no están en el conjunto pero todo intervalo alrededor de ellos lo corta.\quad
$\Z'=\vac$: cada entero $n$ tiene el abierto $(n-\tfrac12,n+\tfrac12)$, que no contiene otro entero; todos los puntos
de $\Z$ son aislados.\quad $\{\tfrac1n:n\ge1\}'=\{0\}$.\quad $\Q'=\R$.
\end{ejemplo}

\begin{metodo}[Método: calcular $A'$ en un espacio finito]
Por la sección~\ref{sec:minimo}, basta probar con el abierto mínimo:
\[x\in A'\iff (U_x\sm\{x\})\cap A\neq\vac .\]
\begin{enumerate}[leftmargin=1.6em]
\item Hacer la tabla de abiertos mínimos $U_x$ (se hace una sola vez por topología y sirve para todo).
\item Para cada $x\in X$: quitar $x$ de $U_x$ e intersecar con $A$.
\item $x\in A'$ exactamente cuando esa intersección no es vacía.
\end{enumerate}
\emph{Atajo:} si $\{x\}$ es abierto, entonces $U_x\sm\{x\}=\vac$ y $x$ no es punto de acumulación de \emph{ningún}
conjunto.
\end{metodo}
La equivalencia es válida porque $U_x$ está contenido en todo abierto que contiene a $x$: si $U_x\sm\{x\}$ ya corta a
$A$, todos los abiertos más grandes también; y si $U_x$ falla, es el testigo de que $x\notin A'$.

\begin{ejemplo}[Ejemplo: $G=\s{b,e}$ en $\tau_1$ y en $\tau_4$ (Parte I)]
\begin{center}\small\renewcommand{\arraystretch}{1.15}
\begin{tabular}{@{}c|ll|ll@{}}
 & \multicolumn{2}{c|}{en $\tau_1$} & \multicolumn{2}{c}{en $\tau_4$}\\
$x$ & $U_x$ & $(U_x\sm\{x\})\cap G$ & $U_x$ & $(U_x\sm\{x\})\cap G$\\ \hline
$a$ & $\s{a}$ & $\vac$ & $\s{a}$ & $\vac$\\
$b$ & $\s{b,c,d,e}$ & $\s{e}$ & $\s{b,e}$ & $\s{e}$\\
$c$ & $\s{c,d}$ & $\vac$ & $\s{b,c,d,e}$ & $\s{b,e}$\\
$d$ & $\s{c,d}$ & $\vac$ & $\s{b,c,d,e}$ & $\s{b,e}$\\
$e$ & $\s{b,c,d,e}$ & $\s{b}$ & $\s{b,e}$ & $\s{b}$
\end{tabular}\end{center}
En $\tau_1$: $G'=\s{b,e}$. En $\tau_4$: $G'=\s{b,c,d,e}$. El mismo conjunto $G$ tiene derivados distintos según la
topología: el derivado no es una propiedad del conjunto solo, sino del conjunto \emph{dentro del espacio}.
\end{ejemplo}

\begin{ejemplo}[Ejemplo: casos extremos]
\begin{itemize}
\item \textbf{Topología discreta:} todo punto es abierto, así que $A'=\vac$ para todo $A$.
\item \textbf{Topología indiscreta:} el único abierto que contiene a $x$ es $X$. Si $A$ tiene al menos dos puntos,
$A'=X$; si $A=\{p\}$, entonces $A'=X\sm\{p\}$; y $\vac'=\vac$.
\item \textbf{Parte II, $\tau_2$, $L=\s{1,4}$:} $L'=\s{1,6}$. El punto $4\in L$ no es de acumulación de $L$ (su abierto
mínimo $\s{2,4,6}$ no contiene otro punto de $L$), y $6\notin L$ sí lo es.
\end{itemize}
\end{ejemplo}

\subsection{Propiedades}
\begin{teorema}[Propiedades del derivado]
Para $A,B\subseteq X$:
\begin{enumerate}[label=(D\arabic*),leftmargin=2.6em]
\item $\vac'=\vac$.
\item Si $A\subseteq B$, entonces $A'\subseteq B'$.
\item $(A\cup B)'=A'\cup B'$.
\item $A$ es cerrado $\iff A'\subseteq A$.
\end{enumerate}
\end{teorema}
\emph{Demostración de (D4).} ($\Rightarrow$) Sea $A$ cerrado y $x\notin A$. Entonces $X\sm A$ es un abierto que
contiene a $x$ y no corta a $A$; luego $x\notin A'$. Así $A'\subseteq A$. ($\Leftarrow$) Sea $A'\subseteq A$ y
$x\notin A$. Como $x\notin A'$, hay un abierto $U_x\ni x$ con $(U_x\sm\{x\})\cap A=\vac$; y como $x\notin A$,
$U_x\cap A=\vac$, es decir, $U_x\subseteq X\sm A$. Por (N4), $X\sm A$ es abierto. \qed

\emph{Demostración de (D3).} «$\supseteq$» sale de (D2). «$\subseteq$»: si $x\notin A'\cup B'$, hay abiertos $U,V$ que
contienen a $x$ con $(U\sm\{x\})\cap A=\vac$ y $(V\sm\{x\})\cap B=\vac$; entonces $U\cap V$ es un abierto que contiene a
$x$ y $((U\cap V)\sm\{x\})\cap(A\cup B)=\vac$, luego $x\notin(A\cup B)'$. \qed

En el taller, (D3) se usó como comprobación: en $\tau_2$ de la Parte~II, $(L\cup M)'=L'\cup M'=\s{1,6}\cup\s{1,4,6}=\s{1,4,6}$.

% =====================================================================
\section{Clausura}
\subsection{Definición y caracterizaciones}
\begin{definicion}[Definición: clausura]
La \textbf{clausura} (o adherencia) de $A\subseteq X$ es la intersección de todos los cerrados que contienen a $A$:
\[\cl{A}=\bigcap\{F:\ F\text{ cerrado},\ A\subseteq F\}.\]
Es un cerrado (C2), contiene a $A$, y es el \emph{menor} cerrado que contiene a $A$.
\end{definicion}
\begin{teorema}[Tres formas equivalentes de calcular la clausura]
Para $A\subseteq X$ y $x\in X$:
\begin{enumerate}[label=(\alph*),leftmargin=1.9em]
\item $x\in\cl{A}\iff$ todo abierto que contiene a $x$ corta a $A$.
\item $\cl{A}=A\cup A'$.
\item En un espacio finito: $x\in\cl{A}\iff U_x\cap A\neq\vac$.
\end{enumerate}
\end{teorema}
\emph{Demostración.} (a) Por contrarrecíproco: $x\notin\cl{A}$ $\iff$ existe un cerrado $F\supseteq A$ con $x\notin F$
$\iff$ existe un abierto $U=X\sm F$ con $x\in U$ y $U\cap A=\vac$.
(b) Si $x\in A$ o $x\in A'$, todo abierto que contiene a $x$ corta a $A$, así que $x\in\cl{A}$ por (a).
Recíprocamente, si $x\in\cl{A}$ y $x\notin A$, todo abierto que contiene a $x$ corta a $A$ en un punto que
necesariamente es distinto de $x$; luego $x\in A'$. (c) Es (a) aplicado al abierto mínimo. \qed

\begin{metodo}[Método: calcular $\cl{A}$]
\textbf{Vía cerrados (la más rápida si ya se tiene $\cotau$):} recorrer la lista de cerrados y quedarse con el más
pequeño que contiene a $A$. Si $A$ mismo aparece en la lista, $\cl{A}=A$.\\[2pt]
\textbf{Vía derivado:} calcular $A'$ y unir: $\cl{A}=A\cup A'$.\\[2pt]
\textbf{Vía abiertos mínimos:} $\cl{A}=\{x: U_x\cap A\neq\vac\}$.\\[2pt]
Hacerlo por dos vías es la mejor comprobación en un examen.
\end{metodo}

\begin{ejemplo}[Ejemplo: Parte I, $\tau_1$, con $\co{\tau_1}=\{\vac,X,\s{a},\s{b,e},\s{a,b,e},\s{b,c,d,e}\}$]
\begin{itemize}
\item $\cl{\s{a}}=\s{a}$: ya es cerrado.
\item $\cl{\s{b}}=\s{b,e}$: los cerrados que contienen a $b$ son $\s{b,e}$, $\s{a,b,e}$, $\s{b,c,d,e}$, $X$; el menor es
$\s{b,e}$. Por la otra vía: $\s{b}'=\s{e}$, y $\s{b}\cup\s{e}=\s{b,e}$.
\item $\cl{\s{c,d}}=\s{b,c,d,e}$. Vía abiertos mínimos: $U_b=\s{b,c,d,e}$ y $U_e=\s{b,c,d,e}$ cortan a $\s{c,d}$;
$U_a=\s{a}$ no.
\end{itemize}
\end{ejemplo}

\subsection{Propiedades (axiomas de Kuratowski)}
\begin{teorema}[Propiedades de la clausura]
\begin{enumerate}[label=(K\arabic*),leftmargin=2.6em]
\item $\cl{\vac}=\vac$.
\item $A\subseteq\cl{A}$.
\item $\cl{\cl{A}}=\cl{A}$.
\item $\cl{A\cup B}=\cl{A}\cup\cl{B}$.
\end{enumerate}
Además: $A$ es cerrado $\iff A=\cl{A}$; y $A\subseteq B\Rightarrow\cl{A}\subseteq\cl{B}$.
\end{teorema}
Estas cuatro propiedades caracterizan a la clausura: un operador que las cumpla define una topología (los cerrados
son los conjuntos con $\cl{A}=A$). Son los \emph{axiomas de clausura de Kuratowski}\fuente{15}.

\begin{ojo}[Con la intersección solo hay una inclusión]
$\cl{A\cap B}\subseteq\cl{A}\cap\cl{B}$, y puede ser estricta. En $\R$: $A=(0,1)$, $B=(1,2)$. Entonces
$A\cap B=\vac$ y $\cl{A\cap B}=\vac$, pero $\cl{A}\cap\cl{B}=[0,1]\cap[1,2]=\{1\}$.
\end{ojo}

% =====================================================================
\section{Interior, exterior y frontera}
\subsection{Definiciones}
\begin{definicion}[Definición: interior, exterior, frontera]
Sea $A\subseteq X$.
\begin{itemize}
\item El \textbf{interior} de $A$ es la unión de todos los abiertos contenidos en $A$:
$\inte{A}=\Int(A)=\bigcup\{U\in\tau: U\subseteq A\}$. Es el \emph{mayor} abierto contenido en $A$. Sus puntos se llaman
puntos interiores: $x\in\inte{A}\iff$ existe un abierto $U$ con $x\in U\subseteq A$ $\iff A\in\N_x$.
\item El \textbf{exterior} de $A$ es el interior de su complemento: $\Ext(A)=\Int(X\sm A)$.
\item La \textbf{frontera} de $A$ es el conjunto de puntos que no son interiores ni exteriores:
$\Fr(A)=X\sm(\inte{A}\cup\Ext(A))$.
\end{itemize}
\end{definicion}
\begin{center}
\begin{tikzpicture}[x=1cm,y=1cm,font=\small]
\fill[tinta!6] (-4.6,-1.7) rectangle (4.6,1.7);
\draw[black!50] (-4.6,-1.7) rectangle (4.6,1.7);
\fill[verde!22] (-0.4,0) ellipse (3.0 and 1.2);
\draw[rojo,line width=1.6pt] (-0.4,0) ellipse (3.0 and 1.2);
\node at (-0.4,0.25) {$\inte{A}$};
\node at (-0.4,-0.3) {\scriptsize puntos con un abierto dentro de $A$};
\node[rojo] at (3.7,1.3) {$\Fr(A)$};
\draw[rojo,-{Stealth[length=2mm]}] (3.35,1.15) -- (2.2,0.62);
\node at (-3.8,-1.35) {$\Ext(A)$};
\node[anchor=west] at (4.8,0.9) {$X=\inte{A}\ \sqcup\ \Fr(A)\ \sqcup\ \Ext(A)$};
\node[anchor=west] at (4.8,0.2) {$\cl{A}=\inte{A}\ \sqcup\ \Fr(A)$};
\node[anchor=west] at (4.8,-0.5) {$\Ext(A)=X\sm\cl{A}$};
\end{tikzpicture}
\end{center}
El dibujo es la imagen mental correcta en $\R^2$; en un espacio finito la frontera no es «un borde delgado» —puede
ser casi todo el espacio—, pero las fórmulas valen igual.

\begin{teorema}[Fórmulas que relacionan interior, clausura y frontera]
\begin{enumerate}[label=(F\arabic*),leftmargin=2.6em]
\item \textbf{Dualidad:} $\Int(A)=X\sm\cl{X\sm A}$ \quad y \quad $\cl{A}=X\sm\Int(X\sm A)=X\sm\Ext(A)$.
\item $\Fr(A)=\cl{A}\sm\inte{A}=\cl{A}\cap\cl{X\sm A}$. En particular $\Fr(A)$ es cerrado y $\Fr(A)=\Fr(X\sm A)$.
\item $x\in\Fr(A)\iff$ todo abierto que contiene a $x$ corta a $A$ y a $X\sm A$.
\item $A$ es abierto $\iff A=\inte{A}\iff A\cap\Fr(A)=\vac$.\quad $A$ es cerrado $\iff\Fr(A)\subseteq A$.\quad
$A$ es abierto y cerrado $\iff\Fr(A)=\vac$.
\item $\Int(A\cap B)=\Int(A)\cap\Int(B)$; \quad $\Int(A)\cup\Int(B)\subseteq\Int(A\cup B)$, posiblemente estricta.
\item En un espacio finito: $x\in\inte{A}\iff U_x\subseteq A$.
\end{enumerate}
\end{teorema}
\emph{Demostración de (F1).} $x\in\Int(A)$ $\iff$ existe un abierto $U\ni x$ con $U\subseteq A$ $\iff$ existe un abierto
$U\ni x$ con $U\cap(X\sm A)=\vac$ $\iff x\notin\cl{X\sm A}$. La segunda fórmula es la primera aplicada a $X\sm A$. \qed

\emph{Demostración de (F2).} Por definición y por (F1), $\Fr(A)=X\sm\inte{A}\ \cap\ X\sm\Ext(A)=\cl{X\sm A}\cap\cl{A}$; y
$\cl{A}\cap(X\sm\inte{A})=\cl{A}\sm\inte{A}$. \qed

\emph{Contraejemplo para (F5).} En $\R$: $A=[0,1]$, $B=[1,2]$. $\Int(A)\cup\Int(B)=(0,1)\cup(1,2)$ pero
$\Int(A\cup B)=(0,2)$.

\begin{metodo}[Método: «anatomía completa» de un conjunto $A$]
\begin{enumerate}[leftmargin=1.6em]
\item \textbf{Interior:} listar los abiertos contenidos en $A$ y unirlos (o tomar el mayor).
\item \textbf{Clausura:} el menor cerrado que contiene a $A$.
\item \textbf{Frontera:} $\cl{A}\sm\inte{A}$.
\item \textbf{Exterior:} $X\sm\cl{A}$.
\item \textbf{Comprobar} que interior, frontera y exterior son disjuntos y que su unión es $X$.
\end{enumerate}
\end{metodo}

\begin{ejemplo}[Ejemplo: $J=\s{a,b,c}$ en $\tau_1$ (Parte I, incisos g y h)]
Interior: los abiertos contenidos en $J$ son $\vac$ y $\s{a}$, así que $\inte{J}=\s{a}$. Clausura: ningún cerrado
propio contiene a $\s{a,b,c}$, así que $\cl{J}=X$. Frontera: $\Fr(J)=X\sm\s{a}=\s{b,c,d,e}$. Exterior:
$X\sm\cl{J}=\vac$.

Aquí la frontera tiene cuatro de los cinco puntos del espacio, y contiene puntos de $J$ ($b$, $c$) y puntos que no
están en $J$ ($d$, $e$). Verificación con (F3) para $d$: su abierto mínimo $\s{c,d}$ corta a $J$ (en $c$) y a
$X\sm J$ (en $d$).
\end{ejemplo}
\begin{ejemplo}[Ejemplo: $L\cup M=\s{1,2,4,6}$ en $\tau_2$ (Parte II, inciso g)]
Interior $\s{2,4,6}$; clausura $\s{1,2,4,6}$ (es cerrado); frontera $\s{1}$; exterior $\s{3}$. Partición:
$\s{2,4,6}\sqcup\s{1}\sqcup\s{3}=X$.
\end{ejemplo}
\begin{ejemplo}[Ejemplos en $\R$]
$A=[0,1)$: interior $(0,1)$, clausura $[0,1]$, frontera $\{0,1\}$, exterior $(-\infty,0)\cup(1,\infty)$.\quad
$A=\Q$: interior $\vac$, clausura $\R$, frontera $\R$, exterior $\vac$.
\end{ejemplo}

\begin{ojo}[Sobre el símbolo $\sigma(J)$ del taller]
Los textos consultados escriben la frontera como $\Fr(A)$, $\partial A$, $b(A)$ o $\mathrm{Bd}\,A$; no se encontró
ninguno que use $\sigma(A)$. En la solución se interpretó $\sigma$ como frontera por el orden de los incisos, pero es
una suposición: hay que confirmarla en los apuntes de clase. Para el parcial, lo seguro es saber calcular los cuatro
objetos (interior, exterior, frontera, clausura) y el derivado.
\end{ojo}

% =====================================================================
\section{Conjuntos densos}
\begin{definicion}[Definición: denso]
$D\subseteq X$ es \textbf{denso} en $X$ si $\cl{D}=X$.
\end{definicion}
\begin{teorema}[Equivalencias]
Son equivalentes: (a) $D$ es denso; (b) todo abierto no vacío corta a $D$; (c) $\Ext(D)=\vac$; (d) el único cerrado
que contiene a $D$ es $X$; (e) en un espacio finito: $D\cap U_x\neq\vac$ para todo $x\in X$.
\end{teorema}
\emph{Demostración de (a)$\iff$(b).} $\cl{D}=X$ $\iff$ todo $x\in X$ cumple que cada abierto que lo contiene corta a
$D$ $\iff$ todo abierto no vacío corta a $D$ (un abierto no vacío contiene a alguno de sus puntos). \qed

\begin{metodo}[Método: ¿cuáles abiertos son densos? (inciso I.e)]
Calcular la clausura de cada abierto; son densos los que dan $X$. Atajo con (b): un abierto $U$ es denso si y solo si
corta a todos los demás abiertos no vacíos.
\end{metodo}
\begin{ejemplo}[Ejemplo: Parte I]
En $\tau_1$: $\cl{\s{a}}=\s{a}$, $\cl{\s{c,d}}=\s{b,c,d,e}$, $\cl{\s{a,c,d}}=X$, $\cl{\s{b,c,d,e}}=\s{b,c,d,e}$. Densos:
$\s{a,c,d}$ y $X$. Con el atajo: $\s{c,d}$ no es denso porque no corta al abierto $\s{a}$; $\s{a,c,d}$ corta a
$\s{a}$, $\s{c,d}$ y $\s{b,c,d,e}$.

En $\tau_4$ los densos son $\s{a,b,e}$ y $X$.
\end{ejemplo}
\begin{ejemplo}[Otros ejemplos]
$\Q$ es denso en $\R$. En la topología del punto incluido (III.7), $\s{0}$ es denso: todo abierto no vacío contiene a
$0$. En la topología discreta el único denso es $X$; en la indiscreta, todo conjunto no vacío es denso.
\end{ejemplo}

\section{Topología relativa (subespacios)}
\subsection{Definición}
Si $H$ es un subconjunto de un espacio $X$, hay una manera natural de convertir a $H$ en espacio: mirar los abiertos
de $X$ «a través de la ventana $H$».
\begin{definicion}[Definición: topología relativa]
Sea $(X,\tau)$ un espacio y $H\subseteq X$. La \textbf{topología relativa} (inducida, o de subespacio) sobre $H$ es
\[\tau_H=\{U\cap H:\ U\in\tau\}.\]
$(H,\tau_H)$ se llama \textbf{subespacio} de $X$.
\end{definicion}
\emph{$\tau_H$ es una topología sobre $H$.} (T1) $\vac=\vac\cap H$ y $H=X\cap H$. (T2)
$\bigcup_i(U_i\cap H)=(\bigcup_iU_i)\cap H$. (T3) $(U\cap H)\cap(V\cap H)=(U\cap V)\cap H$. En los tres casos lo que
queda dentro del paréntesis es un abierto de $X$. \qed

\begin{metodo}[Método: calcular $\tau_H$]
Intersecar \emph{cada} abierto de $\tau$ con $H$, anotar los resultados y eliminar los repetidos. La lista resultante
tiene a lo sumo tantos elementos como $\tau$.
\end{metodo}
\begin{ejemplo}[Ejemplo: Parte I(d) y Parte II(d)]
\textbf{I(d).} $H=\s{a,c,d}$ en $\tau_1$: las intersecciones son $H,\vac,\s{a},\s{c,d},H,\s{c,d}$. Sin repetidos:
$\tau_H=\{H,\vac,\s{a},\s{c,d}\}$. Con $\tau_4$ resulta la \emph{misma} $\tau_H$: dos topologías distintas sobre $X$
pueden inducir la misma topología en un subconjunto.

\textbf{II(d).} $P=\s{1,3,4,6}$ en $\tau_2=\{X,\vac,\s{2},\s{3},\s{2,3},\s{2,4,6},\s{2,3,4,6}\}$: intersecciones
$P,\vac,\vac,\s{3},\s{3},\s{4,6},\s{3,4,6}$. Luego $\tau_P=\{P,\vac,\s{3},\s{4,6},\s{3,4,6}\}$.
\end{ejemplo}

\subsection{Lo que hay que saber de los subespacios}
\begin{teorema}[Propiedades]
Sea $H\subseteq X$ con la topología relativa.
\begin{enumerate}[label=(\alph*),leftmargin=1.9em]
\item Los cerrados de $H$ son exactamente los conjuntos $F\cap H$ con $F$ cerrado en $X$.
\item Si $\B$ es base de $\tau$, entonces $\{B\cap H:B\in\B\}$ es base de $\tau_H$.
\item Para $A\subseteq H$: la clausura de $A$ en $H$ es $\cl{A}\cap H$ (con $\cl{A}$ la clausura en $X$).
\item Si $H$ es abierto en $X$: un subconjunto de $H$ es abierto en $H$ $\iff$ es abierto en $X$.
\end{enumerate}
\end{teorema}
\emph{Demostración de (a).} $C\subseteq H$ es cerrado en $H$ $\iff H\sm C=U\cap H$ para algún $U\in\tau$
$\iff C=H\sm(U\cap H)=(X\sm U)\cap H$. \qed

\begin{ojo}[Abierto en $H$ no implica abierto en $X$]
«Ser abierto» depende del espacio en que se mire. En II(d), $\s{4,6}$ es abierto en $P$ pero no es abierto en $X$
(no está en $\tau_2$). En $\R$: $[0,1)$ es abierto en el subespacio $H=[0,2]$, porque $[0,1)=(-1,1)\cap[0,2]$, y no
es abierto en $\R$. Lo mismo pasa con «cerrado». En un examen conviene escribir siempre «abierto \emph{en} \dots».
\end{ojo}

% =====================================================================
\section{Comparación de topologías}\label{sec:comp}
\begin{definicion}[Definición: más fina, más gruesa, comparables]
Sean $\tau$ y $\tau'$ dos topologías sobre el mismo conjunto $X$. Si $\tau\subseteq\tau'$ se dice que $\tau'$ es
\textbf{más fina} (o más fuerte) que $\tau$, y que $\tau$ es \textbf{más gruesa} (o más débil) que $\tau'$. Las dos
topologías son \textbf{comparables} si $\tau\subseteq\tau'$ o $\tau'\subseteq\tau$.
\end{definicion}
«Más fina» significa «con más abiertos». La imagen de Munkres\fuente{2}: un camión de grava (los abiertos); si se
tritura cada piedra, hay más piedras y más pequeñas: la grava es más fina. Algunos autores de análisis usan «fuerte» y
«débil» con el sentido contrario\fuente{17}; «fina» y «gruesa» no tienen esa ambigüedad.

\begin{teorema}[Hechos básicos]
\begin{enumerate}[label=(\alph*),leftmargin=1.9em]
\item La indiscreta es más gruesa que cualquier topología y la discreta es más fina que cualquiera.
\item La inclusión es un orden parcial entre las topologías de $X$, pero no total: hay pares no comparables.
\item La \textbf{intersección} de dos topologías (o de cualquier familia de topologías) es una topología.
\item La \textbf{unión} de dos topologías puede no ser una topología.
\item $\tau\subseteq\tau'$ $\iff$ la identidad $\mathrm{id}:(X,\tau')\to(X,\tau)$ es continua.
\item Si $\B$ y $\B'$ son bases de $\tau$ y $\tau'$: $\tau\subseteq\tau'$ $\iff$ para todo $B\in\B$ y todo $x\in B$
existe $B'\in\B'$ con $x\in B'\subseteq B$.
\end{enumerate}
\end{teorema}
\emph{Demostración de (c).} Si $\tau=\bigcap_j\tau_j$: $\vac$ y $X$ están en todas; si una familia de conjuntos está en
todas las $\tau_j$, su unión también, y lo mismo la intersección de dos. \qed

\emph{Contraejemplo para (d).} En $X=\s{a,b,c}$: $\tau=\{\vac,\s{a},X\}$ y $\tau'=\{\vac,\s{b},X\}$. La unión
$\{\vac,\s{a},\s{b},X\}$ no contiene a $\s{a}\cup\s{b}=\s{a,b}$. La menor topología que contiene a ambas es la
\emph{generada} por la unión: $\{\vac,\s{a},\s{b},\s{a,b},X\}$.

\begin{metodo}[Método: decidir si dos topologías son comparables]
Mirar si todos los abiertos de una están en la otra. Si $\tau$ tiene un abierto que no está en $\tau'$ \emph{y}
$\tau'$ tiene uno que no está en $\tau$, no son comparables: esos dos conjuntos son el testigo. Atajo: dos topologías
distintas con el mismo número (finito) de abiertos nunca son comparables.
\end{metodo}

\begin{ejemplo}[Ejemplo: las 29 topologías sobre $\s{a,b,c}$ (problema III.4)]
Por inclusión forman un orden parcial con mínimo (indiscreta) y máximo (discreta). De los $\binom{29}{2}=\paresTotal$
pares, $\paresComp$ son comparables y $\paresNoComp$ no. El diagrama de Hasse muestra el orden: se sube por una línea
cuando una topología está contenida en otra sin ninguna intermedia. La numeración es la de la lista de la solución.
\begin{center}\resizebox{0.92\linewidth}{!}{\begin{tikzpicture}[x=1cm,y=1cm,every node/.style={font=\scriptsize}]
\draw[hasse] (0.00,0.00) -- (-7.75,2.30);
\draw[hasse] (0.00,0.00) -- (-4.65,2.30);
\draw[hasse] (0.00,0.00) -- (-1.55,2.30);
\draw[hasse] (0.00,0.00) -- (1.55,2.30);
\draw[hasse] (0.00,0.00) -- (4.65,2.30);
\draw[hasse] (0.00,0.00) -- (7.75,2.30);
\draw[hasse] (-7.75,2.30) -- (-7.75,4.90);
\draw[hasse] (-7.75,2.30) -- (-5.81,4.90);
\draw[hasse] (-7.75,2.30) -- (-3.88,4.90);
\draw[hasse] (-4.65,2.30) -- (-1.94,4.90);
\draw[hasse] (-4.65,2.30) -- (0.00,4.90);
\draw[hasse] (-4.65,2.30) -- (1.94,4.90);
\draw[hasse] (-1.55,2.30) -- (3.88,4.90);
\draw[hasse] (-1.55,2.30) -- (5.81,4.90);
\draw[hasse] (-1.55,2.30) -- (7.75,4.90);
\draw[hasse] (1.55,2.30) -- (-7.75,4.90);
\draw[hasse] (1.55,2.30) -- (-1.94,4.90);
\draw[hasse] (1.55,2.30) -- (3.88,4.90);
\draw[hasse] (4.65,2.30) -- (-5.81,4.90);
\draw[hasse] (4.65,2.30) -- (0.00,4.90);
\draw[hasse] (4.65,2.30) -- (5.81,4.90);
\draw[hasse] (7.75,2.30) -- (-3.88,4.90);
\draw[hasse] (7.75,2.30) -- (1.94,4.90);
\draw[hasse] (7.75,2.30) -- (7.75,4.90);
\draw[hasse] (-7.75,4.90) -- (-7.75,7.50);
\draw[hasse] (-7.75,4.90) -- (-1.55,7.50);
\draw[hasse] (-5.81,4.90) -- (-4.65,7.50);
\draw[hasse] (-5.81,4.90) -- (-1.55,7.50);
\draw[hasse] (-3.88,4.90) -- (-4.65,10.10);
\draw[hasse] (-3.88,4.90) -- (1.55,10.10);
\draw[hasse] (-1.94,4.90) -- (-7.75,7.50);
\draw[hasse] (-1.94,4.90) -- (4.65,7.50);
\draw[hasse] (0.00,4.90) -- (-7.75,10.10);
\draw[hasse] (0.00,4.90) -- (7.75,10.10);
\draw[hasse] (1.94,4.90) -- (1.55,7.50);
\draw[hasse] (1.94,4.90) -- (4.65,7.50);
\draw[hasse] (3.88,4.90) -- (-1.55,10.10);
\draw[hasse] (3.88,4.90) -- (4.65,10.10);
\draw[hasse] (5.81,4.90) -- (-4.65,7.50);
\draw[hasse] (5.81,4.90) -- (7.75,7.50);
\draw[hasse] (7.75,4.90) -- (1.55,7.50);
\draw[hasse] (7.75,4.90) -- (7.75,7.50);
\draw[hasse] (-7.75,7.50) -- (-7.75,10.10);
\draw[hasse] (-7.75,7.50) -- (-4.65,10.10);
\draw[hasse] (-4.65,7.50) -- (-1.55,10.10);
\draw[hasse] (-4.65,7.50) -- (1.55,10.10);
\draw[hasse] (-1.55,7.50) -- (-7.75,10.10);
\draw[hasse] (-1.55,7.50) -- (-1.55,10.10);
\draw[hasse] (1.55,7.50) -- (4.65,10.10);
\draw[hasse] (1.55,7.50) -- (7.75,10.10);
\draw[hasse] (4.65,7.50) -- (-4.65,10.10);
\draw[hasse] (4.65,7.50) -- (4.65,10.10);
\draw[hasse] (7.75,7.50) -- (1.55,10.10);
\draw[hasse] (7.75,7.50) -- (7.75,10.10);
\draw[hasse] (-7.75,10.10) -- (0.00,12.40);
\draw[hasse] (-4.65,10.10) -- (0.00,12.40);
\draw[hasse] (-1.55,10.10) -- (0.00,12.40);
\draw[hasse] (1.55,10.10) -- (0.00,12.40);
\draw[hasse] (4.65,10.10) -- (0.00,12.40);
\draw[hasse] (7.75,10.10) -- (0.00,12.40);
\node[nodoh] at (0.00,0.00) {$\tau_{1}$};
\node[nodoh] at (-7.75,2.30) {$\tau_{2}$};
\node[nodoh] at (-4.65,2.30) {$\tau_{3}$};
\node[nodoh] at (-1.55,2.30) {$\tau_{4}$};
\node[nodoh] at (1.55,2.30) {$\tau_{5}$};
\node[nodoh] at (4.65,2.30) {$\tau_{6}$};
\node[nodoh] at (7.75,2.30) {$\tau_{7}$};
\node[nodoh] at (-7.75,4.90) {$\tau_{8}$};
\node[nodoh] at (-5.81,4.90) {$\tau_{9}$};
\node[nodoh] at (-3.88,4.90) {$\tau_{10}$};
\node[nodoh] at (-1.94,4.90) {$\tau_{11}$};
\node[nodoh] at (0.00,4.90) {$\tau_{12}$};
\node[nodoh] at (1.94,4.90) {$\tau_{13}$};
\node[nodoh] at (3.88,4.90) {$\tau_{14}$};
\node[nodoh] at (5.81,4.90) {$\tau_{15}$};
\node[nodoh] at (7.75,4.90) {$\tau_{16}$};
\node[nodoh] at (-7.75,7.50) {$\tau_{17}$};
\node[nodoh] at (-4.65,7.50) {$\tau_{18}$};
\node[nodoh] at (-1.55,7.50) {$\tau_{19}$};
\node[nodoh] at (1.55,7.50) {$\tau_{20}$};
\node[nodoh] at (4.65,7.50) {$\tau_{21}$};
\node[nodoh] at (7.75,7.50) {$\tau_{22}$};
\node[nodoh] at (-7.75,10.10) {$\tau_{23}$};
\node[nodoh] at (-4.65,10.10) {$\tau_{24}$};
\node[nodoh] at (-1.55,10.10) {$\tau_{25}$};
\node[nodoh] at (1.55,10.10) {$\tau_{26}$};
\node[nodoh] at (4.65,10.10) {$\tau_{27}$};
\node[nodoh] at (7.75,10.10) {$\tau_{28}$};
\node[nodoh] at (0.00,12.40) {$\tau_{29}$};
\node[anchor=east,font=\scriptsize\itshape,text=black!60] at (-8.45,0.00) {2 abiertos};
\node[anchor=east,font=\scriptsize\itshape,text=black!60] at (-8.45,2.30) {3 abiertos};
\node[anchor=east,font=\scriptsize\itshape,text=black!60] at (-8.45,4.90) {4 abiertos};
\node[anchor=east,font=\scriptsize\itshape,text=black!60] at (-8.45,7.50) {5 abiertos};
\node[anchor=east,font=\scriptsize\itshape,text=black!60] at (-8.45,10.10) {6 abiertos};
\node[anchor=east,font=\scriptsize\itshape,text=black!60] at (-8.45,12.40) {8 abiertos};
\end{tikzpicture}
}\end{center}
Lecturas: $\tau_2=\{\vac,\s{a},X\}\subset\tau_8=\{\vac,\s{a},\s{a,b},X\}\subset\tau_{17}=\{\vac,\s{a},\s{b},\s{a,b},X\}$
es una cadena de topologías cada vez más finas. En cambio $\tau_2$ y $\tau_3=\{\vac,\s{b},X\}$ no son comparables.
\end{ejemplo}

% =====================================================================
\section{Bases y sub-bases}\label{sec:bases}
\subsection{Para qué sirven}
Describir una topología listando \emph{todos} sus abiertos solo es posible en ejemplos pequeños. En $\R$ nadie lista
los abiertos: se dice «las uniones de intervalos abiertos». Los intervalos son una \emph{base}: un catálogo reducido de
abiertos con el cual se arman todos los demás por unión. Una \emph{sub-base} es un catálogo aún más reducido, con el
que se arma la base por intersecciones finitas.
\[\text{sub-base }\Sb\ \xrightarrow{\ \text{intersecciones finitas}\ }\ \text{base }\B\
\xrightarrow{\ \text{uniones arbitrarias}\ }\ \text{topología }\tau\]

\subsection{Bases}
\begin{definicion}[Definición: base de una topología]
Sea $(X,\tau)$ un espacio. Una familia $\B\subseteq\tau$ es una \textbf{base} de $\tau$ si todo abierto es unión de
miembros de $\B$. Equivalentemente: para todo abierto $U$ y todo $x\in U$ existe $B\in\B$ con $x\in B\subseteq U$.
\end{definicion}
(El vacío es la unión de la subfamilia vacía, así que $\vac$ no necesita estar en la base.)

\begin{ejemplo}[Ejemplos de bases]
\begin{itemize}
\item Los intervalos abiertos $(a,b)$ son base de la topología usual de $\R$; también lo son los de extremos
racionales.
\item $\{\{x\}:x\in X\}$ es base de la topología discreta. $\{X\}$ es base de la indiscreta.
\item $\tau$ misma siempre es base de $\tau$.
\item \textbf{En un espacio finito, los abiertos mínimos $\{U_x:x\in X\}$ forman una base}, y es la más pequeña: todo
abierto $U$ cumple $U=\bigcup_{x\in U}U_x$, y cualquier base debe contener a cada $U_x$. Para $\tau_1$ de la Parte~I:
$\B=\{\s{a},\s{c,d},\s{b,c,d,e}\}$.
\item En el problema III.6, los bloques $\s{2n-1,2n}$ son una base; en el III.7, los conjuntos $\s{0,x}$.
\end{itemize}
\end{ejemplo}

No toda familia de subconjuntos es base de alguna topología. El criterio es:
\begin{teorema}[¿Cuándo una familia $\B$ es base de \emph{alguna} topología sobre $X$?]
Si y solo si:
\begin{description}[leftmargin=1.4cm,labelwidth=1.0cm,itemsep=0.1em]
\item[(B1)] $\B$ cubre $X$: todo $x\in X$ pertenece a algún $B\in\B$;
\item[(B2)] si $x\in B_1\cap B_2$ con $B_1,B_2\in\B$, existe $B_3\in\B$ con $x\in B_3\subseteq B_1\cap B_2$.
\end{description}
En ese caso la topología es única: sus abiertos son las uniones de miembros de $\B$.
\end{teorema}
\emph{Idea de la demostración.} (B1) garantiza que $X$ es unión de miembros de $\B$. (B2) dice que $B_1\cap B_2$ es
unión de miembros de $\B$; de ahí, la intersección de dos uniones de básicos es unión de básicos, que es (T3). (T2) es
automático. \qed

\emph{Ejemplo de familia que no es base.} En el problema III.11, $\C=\{[a,a+1]\}$ cumple (B1) pero no (B2):
$1\in[0,1]\cap[1,2]=\{1\}$ y ningún miembro de $\C$ cabe dentro de $\{1\}$.

\begin{teorema}[Problema III.12: agrandar una base]
Si $\B$ es base de $\tau$ y $\B\subseteq\B_1\subseteq\tau$, entonces $\B_1$ es base de $\tau$.
\end{teorema}
\emph{Demostración.} Los miembros de $\B_1$ son abiertos. Todo abierto es unión de miembros de $\B$, que también son
miembros de $\B_1$. \qed

\subsection{Sub-bases y topología generada}
\begin{definicion}[Definición: sub-base]
Una familia $\Sb\subseteq\tau$ es una \textbf{sub-base} de $\tau$ si las intersecciones finitas de miembros de $\Sb$
forman una base de $\tau$.
\end{definicion}
\begin{teorema}[Topología generada]
Sea $\C$ \emph{cualquier} familia de subconjuntos de $X$. Existe una menor topología sobre $X$ que contiene a $\C$; se
llama \textbf{topología generada} por $\C$ y se denota $\tau(\C)$. Se obtiene así:
\begin{enumerate}[leftmargin=1.6em]
\item $\B$ = todas las intersecciones finitas de miembros de $\C$, junto con $X$;
\item $\tau(\C)$ = todas las uniones de miembros de $\B$ (incluida $\vac$).
\end{enumerate}
Además $\tau(\C)$ es la intersección de todas las topologías que contienen a $\C$, y $\C$ es sub-base de $\tau(\C)$.
\end{teorema}
\emph{Por qué funciona.} La familia $\B$ del paso 1 cubre $X$ y es cerrada bajo intersecciones finitas, luego cumple
(B1) y (B2) y es base de una topología. Cualquier topología que contenga a $\C$ debe contener a las intersecciones
finitas de sus miembros y luego a las uniones de estas: por eso la obtenida es la menor. \qed

\begin{ojo}[¿Y el conjunto $X$?]
Si los miembros de $\C$ no cubren $X$ (como en el problema III.9, donde el punto $e$ no está en ninguno), $X$ no
aparece como intersección de miembros de $\C$. Hay dos maneras equivalentes de resolverlo: (1) convenir que la
intersección de la familia vacía es $X$; (2) exigir en la definición que la sub-base cubra $X$. Los textos difieren en
esto\fuente{16}; con la definición «menor topología que contiene a $\C$» no hay ambigüedad: $X$ siempre está.
\end{ojo}

\begin{metodo}[Método: hallar la topología generada por una familia finita $\C$]
\begin{enumerate}[leftmargin=1.6em]
\item \textbf{Intersecciones:} calcular las intersecciones de dos en dos, de tres en tres, etc., de miembros de $\C$.
Anotar los conjuntos nuevos. Agregar $X$. Eso es $\B$.
\item \textbf{Uniones:} calcular las uniones de miembros de $\B$ hasta que no aparezcan conjuntos nuevos. Agregar
$\vac$. Eso es $\tau(\C)$.
\item \textbf{Verificar} que el resultado es cerrado bajo $\cup$ y $\cap$ (tabla), como pide el enunciado.
\end{enumerate}
El orden importa: primero intersecciones, luego uniones. Al revés puede faltar algún abierto.
\end{metodo}

\begin{ejemplo}[Ejemplo: problema III.8, $\C=\{\s{a,b,c},\s{c,d},\s{d,e}\}$ en $X=\s{a,b,c,d,e}$]
\emph{Intersecciones:} $\s{a,b,c}\cap\s{c,d}=\s{c}$; $\s{c,d}\cap\s{d,e}=\s{d}$; $\s{a,b,c}\cap\s{d,e}=\vac$.
$\B=\{X,\vac,\s{c},\s{d},\s{c,d},\s{d,e},\s{a,b,c}\}$.

\emph{Uniones nuevas:} $\s{c}\cup\s{d,e}=\s{c,d,e}$; $\s{d}\cup\s{a,b,c}=\s{a,b,c,d}$.

$\tau(\C)=\{\vac,\s{c},\s{d},\s{c,d},\s{d,e},\s{c,d,e},\s{a,b,c},\s{a,b,c,d},X\}$: nueve abiertos a partir de tres.
\end{ejemplo}
\begin{ejemplo}[{Ejemplo: problema III.11, $\C=\{[a,a+1]:a\in\R\}$}]
$[x-1,x]\cap[x,x+1]=\{x\}$: todos los puntos son básicos, y la topología generada es la discreta. Una sub-base formada
por conjuntos «grandes» puede generar la topología más fina posible.
\end{ejemplo}
\begin{ejemplo}[Ejemplo: una sub-base de la topología usual de $\R$]
Las semirrectas abiertas $(a,\infty)$ y $(-\infty,b)$. En efecto, $(a,\infty)\cap(-\infty,b)=(a,b)$, y los intervalos
son base.
\end{ejemplo}

\begin{metodo}[Método: encontrar sub-bases de una topología dada (problema III.10)]
Una familia $\Sb\subseteq\tau$ es sub-base de $\tau$ si genera exactamente $\tau$. Maneras de producirlas:
\begin{itemize}
\item $\tau$ misma, o cualquier base de $\tau$ (toda base es sub-base).
\item Partir de los abiertos mínimos y quitar los que se recuperen como intersección de otros abiertos de la familia.
\item Buscar abiertos «grandes» cuyas intersecciones den los pequeños: en la discreta sobre $X$, los complementos de
los puntos.
\end{itemize}
Siempre hay que comprobar al final que al generar no falta ni sobra ningún abierto.
\end{metodo}

\subsection{Por qué importan: continuidad con menos trabajo}
\begin{teorema}[Continuidad sobre una base o sub-base]
Sea $f:X\to Y$ y $\Sb$ una sub-base (o una base) de la topología de $Y$. Entonces $f$ es continua $\iff$
$f^{-1}(S)$ es abierto en $X$ para todo $S\in\Sb$.
\end{teorema}
\emph{Demostración.} La imagen inversa respeta uniones e intersecciones:
$f^{-1}(\bigcup_iV_i)=\bigcup_if^{-1}(V_i)$ y $f^{-1}(V\cap W)=f^{-1}(V)\cap f^{-1}(W)$. Todo abierto de $Y$ es unión
de intersecciones finitas de miembros de $\Sb$; su imagen inversa es unión de intersecciones finitas de abiertos. \qed

Este resultado explica, por ejemplo, por qué en el problema IV.3 la continuidad de $h$ y de $k$ (que llegan a $X$)
se puede decidir con la lista impresa de $\tau_1$ aunque le falte $\s{1,2}$: esa lista es una sub-base de la topología
reparada $\tau_1^{*}$.

\section{Continuidad}\label{sec:cont}
\subsection{De $\varepsilon$--$\delta$ a los abiertos}
En cálculo, $f:\R\to\R$ es continua en $p$ si para todo $\varepsilon>0$ existe $\delta>0$ tal que
$|x-p|<\delta\Rightarrow|f(x)-f(p)|<\varepsilon$. En palabras: \emph{para todo intervalo $V$ alrededor de $f(p)$ existe
un intervalo $U$ alrededor de $p$ con $f(U)\subseteq V$}. Si se cambia «intervalo alrededor de» por «abierto que
contiene a», la frase ya no menciona distancias y tiene sentido en cualquier espacio topológico.

\begin{definicion}[Definición: continuidad en un punto y continuidad global]
Sea $f:(X,\tau_X)\to(Y,\tau_Y)$.
\begin{itemize}
\item $f$ es \textbf{continua en $p\in X$} si para todo abierto $V$ de $Y$ con $f(p)\in V$ existe un abierto $U$ de $X$
con $p\in U$ y $f(U)\subseteq V$.
\item $f$ es \textbf{continua} si la imagen inversa de todo abierto de $Y$ es un abierto de $X$:
\[V\in\tau_Y\ \Longrightarrow\ f^{-1}(V)\in\tau_X .\]
\end{itemize}
\end{definicion}
Recuérdese que $f^{-1}(V)=\{x\in X: f(x)\in V\}$ está definido para \emph{cualquier} función; no supone que $f$ tenga
inversa.

\begin{teorema}[Equivalencias de la continuidad]
Para $f:X\to Y$ son equivalentes:
\begin{enumerate}[label=(\alph*),leftmargin=1.9em]
\item $f$ es continua (imagen inversa de abierto es abierto);
\item la imagen inversa de todo \emph{cerrado} de $Y$ es cerrada en $X$;
\item la imagen inversa de todo miembro de una base (o de una sub-base) de $Y$ es abierta en $X$;
\item $f$ es continua en cada punto de $X$;
\item $f(\cl{A})\subseteq\cl{f(A)}$ para todo $A\subseteq X$.
\end{enumerate}
\end{teorema}
\emph{Demostración de (a)$\iff$(b).} $f^{-1}(Y\sm V)=X\sm f^{-1}(V)$. \qed

\emph{Demostración de (a)$\iff$(d).} ($\Rightarrow$) Dado $p$ y $V\ni f(p)$ abierto, $U=f^{-1}(V)$ es abierto, contiene
a $p$ y $f(U)\subseteq V$. ($\Leftarrow$) Sea $V$ abierto de $Y$ y $p\in f^{-1}(V)$. Por la continuidad en $p$ hay un
abierto $U_p\ni p$ con $f(U_p)\subseteq V$, es decir, $U_p\subseteq f^{-1}(V)$. Entonces $f^{-1}(V)$ es vecindad de cada
uno de sus puntos y, por (N4), es abierto. \qed

\begin{metodo}[Método: decidir si $f$ es continua (espacios finitos)]
\begin{enumerate}[leftmargin=1.6em]
\item Escribir la tabla de valores de $f$ (leerla del diagrama de flechas con cuidado).
\item Para \emph{cada} abierto $V$ de la llegada, calcular $f^{-1}(V)$: los puntos de la salida cuya flecha cae en $V$.
\item Si todas las imágenes inversas son abiertos de la salida, $f$ es continua. Si alguna no lo es, $f$ no es
continua y ese $V$ es el contraejemplo.
\end{enumerate}
Para $\vac$ y para el total no hay nada que revisar: $f^{-1}(\vac)=\vac$ y $f^{-1}(Y)=X$.
\end{metodo}
\begin{metodo}[Método: decidir si $f$ es continua en un punto $p$]
\begin{enumerate}[leftmargin=1.6em]
\item Calcular $f(p)$ y listar los abiertos $V$ de la llegada que contienen a $f(p)$.
\item Para cada uno, ver si $f^{-1}(V)$ es vecindad de $p$ (contiene un abierto que contiene a $p$).
\item En espacios finitos basta una sola prueba: $f$ es continua en $p$ $\iff f(U_p)\subseteq V_{f(p)}$, donde $U_p$ y
$V_{f(p)}$ son los abiertos mínimos de $p$ y de $f(p)$.
\end{enumerate}
\end{metodo}

\begin{ejemplo}[Ejemplo: problema IV.1]
$X=\s{a,b,c,d}$, $\tau=\{X,\vac,\s{a},\s{b},\s{a,b},\s{b,c,d}\}$ y $f(a)=b$, $f(b)=d$, $f(c)=b$, $f(d)=c$.

\emph{Global.} $f^{-1}(\s{b})=\s{a,c}$ (los puntos que van a $b$ son $a$ y $c$), y $\s{a,c}\notin\tau$. No es continua.

\emph{En $d$.} $f(d)=c$; abiertos que contienen a $c$: $\s{b,c,d}$ y $X$. $f^{-1}(\s{b,c,d})=X$, que es vecindad de
$d$. Sí es continua en $d$. Con abiertos mínimos: $U_d=\s{b,c,d}$, $f(U_d)=\s{d,b,c}\subseteq U_{f(d)}=U_c=\s{b,c,d}$.

\emph{En $c$.} $f(c)=b$; $U_c=\s{b,c,d}$, $f(U_c)=\s{b,c,d}\not\subseteq U_b=\s{b}$. No es continua en $c$.

La función es continua en $a$, $b$ y $d$, y discontinua solo en $c$: basta un punto de discontinuidad para que la
función no sea continua.
\end{ejemplo}

\begin{ejemplo}[Ejemplos generales que conviene recordar]
\begin{itemize}
\item Toda función constante es continua (la imagen inversa de un abierto es $\vac$ o $X$).
\item La identidad $(X,\tau)\to(X,\tau)$ es continua.
\item Si $X$ tiene la topología \textbf{discreta}, toda función $X\to Y$ es continua.
\item Si $Y$ tiene la topología \textbf{indiscreta}, toda función $X\to Y$ es continua.
\item La composición de funciones continuas es continua: $(g\circ f)^{-1}(W)=f^{-1}(g^{-1}(W))$.
\item La restricción de una función continua a un subespacio es continua.
\end{itemize}
\end{ejemplo}

\begin{ojo}[Errores frecuentes en continuidad]
\begin{itemize}
\item Revisar imágenes directas $f(U)$ en lugar de imágenes inversas $f^{-1}(V)$. Eso es la definición de función
\emph{abierta}, que es otra cosa.
\item Usar los abiertos del espacio equivocado: las $V$ se toman en la \emph{llegada} y se comprueba que $f^{-1}(V)$
esté en la topología de la \emph{salida}.
\item Leer mal el diagrama: $f^{-1}(\s{b})$ son todos los puntos de los que \emph{sale} una flecha que llega a $b$.
\item Creer que «continua en algunos puntos» implica continua.
\end{itemize}
\end{ojo}

% =====================================================================
\section{Funciones abiertas y funciones cerradas}
\begin{definicion}[Definición]
$f:X\to Y$ es \textbf{abierta} si $f(U)$ es abierto en $Y$ para todo abierto $U$ de $X$; es \textbf{cerrada} si
$f(F)$ es cerrado en $Y$ para todo cerrado $F$ de $X$.
\end{definicion}
Continua, abierta y cerrada son tres propiedades \textbf{independientes}: una función puede tener cualquiera de las
ocho combinaciones. La razón de fondo es que la imagen inversa se comporta bien con todas las operaciones de conjuntos,
mientras que la imagen directa solo respeta las uniones:
\begin{gather*}
f(A\cup B)=f(A)\cup f(B),\qquad f(A\cap B)\subseteq f(A)\cap f(B),\\
f(X\sm A)\ \text{y}\ Y\sm f(A)\ \text{no se relacionan en general}.
\end{gather*}
Por eso «abierta» no equivale a «cerrada», salvo en el caso biyectivo.

\begin{teorema}[Hechos útiles]
\begin{enumerate}[label=(\alph*),leftmargin=1.9em]
\item Para saber si $f$ es abierta basta revisar las imágenes de los miembros de una \emph{base} de $X$ (porque
$f(\bigcup_iB_i)=\bigcup_if(B_i)$).
\item Si $f$ es abierta, entonces $f(X)$ es abierto en $Y$; si es cerrada, $f(X)$ es cerrado en $Y$. (Dos pruebas
rápidas para descartar.)
\item Si $f$ es \textbf{biyectiva}: $f$ abierta $\iff f$ cerrada $\iff f^{-1}$ continua.
\end{enumerate}
\end{teorema}
\emph{Demostración de (c).} Para una biyección, $f(X\sm A)=Y\sm f(A)$, lo que da la primera equivalencia. Y la imagen
inversa de $U$ por la función $f^{-1}:Y\to X$ es precisamente $f(U)$; así, «$f^{-1}$ continua» dice «$f(U)$ abierto para
todo abierto $U$». \qed

\begin{metodo}[Método: decidir si $f$ es abierta / cerrada]
\textbf{Abierta:} para cada abierto $U$ de la salida calcular $f(U)$ (adónde llegan las flechas que salen de $U$) y
ver si está en la topología de la llegada.\\[2pt]
\textbf{Cerrada:} listar los cerrados de la salida y los de la llegada; para cada cerrado $F$ de la salida calcular
$f(F)$ y ver si está en la lista de cerrados de la llegada.\\[2pt]
\textbf{Descartes rápidos:} mirar $f(X)$; mirar los abiertos y cerrados de un solo punto.
\end{metodo}

\begin{ejemplo}[Ejemplo: problema IV.2, $f(x)=x^2$ en $\R$]
$f$ es continua y \textbf{no es abierta}: $f((-1,1))=[0,1)$, que no es abierto porque ningún intervalo alrededor de $0$
cabe en $[0,1)$\fuente{20}. Geométricamente: la parábola «dobla» la recta en $0$, y el punto del doblez queda en el
borde de la imagen.
\end{ejemplo}

\begin{ejemplo}[Ejemplo: problema IV.3, resumen]
Con $X=\s{1,2,3,4}$, $Y=\s{a,b,c,d}$, $\tau_2=\{Y,\vac,\s{c},\s{c,d},\s{a,b,c}\}$ y la $\tau_1$ del enunciado:
\begin{center}\footnotesize\renewcommand{\arraystretch}{1.25}\setlength{\tabcolsep}{4pt}
\begin{tabular}{@{}lcccl@{}}\toprule
función & continua & abierta & cerrada & por qué\\ \midrule
$f=(1\!\mapsto\! a,2\!\mapsto\! a,3\!\mapsto\! d,4\!\mapsto\! c)$ & \no & \no & \no & $f^{-1}(\s{c})=\s{4}$; $f(\s{2})=\s{a}$; $f(\s{1})=\s{a}$\\
$g=(1\!\mapsto\! b,2\!\mapsto\! c,3\!\mapsto\! c,4\!\mapsto\! d)$ & \si & \no & \no & $g(X)=\s{b,c,d}$: ni abierto ni cerrado\\
$h=(a\!\mapsto\! 2,b\!\mapsto\! 3,c\!\mapsto\! 4,d\!\mapsto\! 1)$ & \no & \no & \no & $h^{-1}(\s{2})=\s{a}$; $h(\s{c})=\s{4}$\\
$k=(a\!\mapsto\! 2,b\!\mapsto\! 4,c\!\mapsto\! 3,d\!\mapsto\! 3)$ & \no & \si & \no & $k^{-1}(\s{2})=\s{a}$; $k(Y)=\s{2,3,4}$ no cerrado\\ \bottomrule
\end{tabular}\end{center}
$g$ es continua y no abierta; $k$ es abierta y no continua: el mismo problema contiene los dos contraejemplos que
muestran la independencia de ambas nociones. Nótese el uso del descarte (b): $g(X)$ y $k(Y)$.
\end{ejemplo}

% =====================================================================
\section{Homeomorfismos}
\subsection{Definición e idea}
Dos espacios son «el mismo espacio con los puntos renombrados» cuando hay una correspondencia entre sus puntos que
también hace corresponder sus abiertos. Eso es un homeomorfismo: el isomorfismo de la topología.
\begin{definicion}[Definición: homeomorfismo]
$f:X\to Y$ es un \textbf{homeomorfismo} si es biyectiva, continua y su inversa $f^{-1}:Y\to X$ es continua. Si existe
un homeomorfismo se dice que $X$ y $Y$ son \textbf{homeomorfos} (o topológicamente equivalentes): $X\homeo Y$.
\end{definicion}
\begin{teorema}[Formas equivalentes]
Para una \emph{biyección} $f:X\to Y$ son equivalentes:
\begin{enumerate}[label=(\alph*),leftmargin=1.9em]
\item $f$ es homeomorfismo;
\item $f$ es continua y abierta;
\item $f$ es continua y cerrada;
\item $U$ es abierto en $X$ $\iff$ $f(U)$ es abierto en $Y$; es decir, $U\mapsto f(U)$ es una biyección de $\tau_X$
sobre $\tau_Y$.
\end{enumerate}
\end{teorema}
\begin{teorema}[$\homeo$ es una relación de equivalencia]
Reflexiva: la identidad es homeomorfismo. Simétrica: si $f$ es homeomorfismo, $f^{-1}$ también. Transitiva: la
composición de homeomorfismos es homeomorfismo.
\end{teorema}

\begin{ojo}[Biyectiva y continua no basta]
La inversa debe ser continua también. Ejemplo finito: la identidad de $X=\s{a,b}$ con la topología discreta hacia
$X$ con la indiscreta es biyectiva y continua, pero su inversa no es continua. Ejemplo clásico:
$t\mapsto(\cos t,\sen t)$ de $[0,2\pi)$ sobre la circunferencia es biyectiva y continua, y no es homeomorfismo
(la inversa «rompe» la circunferencia en el punto $(1,0)$)\fuente{22,23}.
\end{ojo}

\subsection{Invariantes topológicos: cómo probar que \emph{no} son homeomorfos}
Una propiedad es un \textbf{invariante topológico} si, cuando la tiene un espacio, la tienen todos los homeomorfos a
él. Si $X$ tiene un invariante y $Y$ no, entonces $X\not\homeo Y$. Para espacios finitos los más útiles son:
\begin{center}\small\renewcommand{\arraystretch}{1.25}
\begin{tabular}{@{}p{6.4cm}p{9.0cm}@{}}\toprule
\textbf{Invariante} & \textbf{Por qué se conserva}\\ \midrule
Número de puntos & Un homeomorfismo es una biyección.\\
Número de abiertos (y de cerrados) & $U\mapsto f(U)$ es biyección entre las topologías.\\
Número de abiertos de cada tamaño & $f(U)$ tiene tantos puntos como $U$.\\
Número de puntos abiertos; de puntos cerrados & Caso particular del anterior.\\
Número de conjuntos abiertos y cerrados a la vez & $f$ conserva abiertos y cerrados.\\
Existencia de dos abiertos no vacíos disjuntos & $f(U\cap V)=f(U)\cap f(V)$ para biyecciones.\\
La forma del orden de inclusión entre abiertos & $U\subseteq V\iff f(U)\subseteq f(V)$.\\ \bottomrule
\end{tabular}\end{center}

\begin{ojo}[El mismo número de abiertos no garantiza homeomorfismo]
En $X=\s{a,b,c}$, $\tau=\{\vac,\s{a},\s{a,b},X\}$ y $\tau'=\{\vac,\s{a},\s{b,c},X\}$ tienen 4 abiertos cada una y no son
homeomorfas: $\tau'$ tiene un abierto propio que también es cerrado ($\s{a}$) y $\tau$ no; o bien, $\tau'$ tiene dos
abiertos no vacíos disjuntos y $\tau$ no. Los invariantes sirven para \emph{negar}; para \emph{afirmar} hay que
construir el homeomorfismo.
\end{ojo}

\begin{metodo}[Método: decidir si $X\homeo Y$ (espacios finitos)]
\begin{enumerate}[leftmargin=1.6em]
\item \textbf{Descartar:} comparar número de puntos, número de abiertos y tamaños de los abiertos. Si algo difiere,
no son homeomorfos y ese es el argumento.
\item \textbf{Construir:} si todo coincide, emparejar los abiertos «por tamaño y posición»: los puntos abiertos con
los puntos abiertos, los pares abiertos con los pares abiertos, etc. Eso dicta adónde debe ir cada punto.
\item \textbf{Verificar:} escribir la tabla $U\mapsto\varphi(U)$ para todos los abiertos de $X$ y comprobar que la
fila de imágenes es exactamente la topología de $Y$.
\end{enumerate}
\end{metodo}

\begin{ejemplo}[Ejemplo: problema IV.4]
$\tau_1=\{X,\vac,\s{a},\s{c,d},\s{a,c,d},\s{b,c,d,e}\}$ y $\tau_2=\{Y,\vac,\s{v},\s{w,z},\s{v,w,z},\s{w,x,y,z}\}$ tienen
abiertos de tamaños $0,1,2,3,4,5$ en ambos casos. Emparejando: $\s{a}\leftrightarrow\s{v}$ obliga $a\mapsto v$;
$\s{c,d}\leftrightarrow\s{w,z}$ obliga $\s{c,d}\mapsto\s{w,z}$; lo que queda, $\s{b,e}\mapsto\s{x,y}$.
Con $\varphi=(a\mapsto v,b\mapsto x,c\mapsto w,d\mapsto z,e\mapsto y)$ la tabla de imágenes recorre $\tau_2$:
$X\homeo Y$.

$\tau_3$ tiene 7 abiertos y las otras 6: $Z\not\homeo X$ y $Z\not\homeo Y$. Y todo espacio es homeomorfo a sí mismo.
\end{ejemplo}
\begin{ejemplo}[Ejemplo: problema IV.3(b)]
$X$ tiene 9 abiertos (10 con la corrección) y $Y$ tiene 5: no pueden ser homeomorfos, a pesar de lo que pide probar
el enunciado. Cuando un enunciado pide demostrar algo falso, la respuesta correcta es el contraargumento, bien
justificado.
\end{ejemplo}
\begin{ejemplo}[Ejemplos en $\R$]
$(-1,1)\homeo\R$ mediante $x\mapsto\dfrac{x}{1-|x|}$ (inversa $y\mapsto\dfrac{y}{1+|y|}$, ambas continuas). Dos
intervalos abiertos cualesquiera son homeomorfos (por una función afín). En cambio $[0,1]\not\homeo(0,1)$: la
compacidad es un invariante que se estudia más adelante en el curso.
\end{ejemplo}

\subsection{Clasificación de los espacios de tres puntos}
Las 29 topologías sobre $\s{a,b,c}$ se agrupan en \textbf{9 clases de homeomorfismo}\fuente{12,13}. Un representante de
cada una, con el número de topologías de la clase:
\begin{center}\small\renewcommand{\arraystretch}{1.2}
\begin{tabular}{@{}llc@{}}\toprule
representante & descripción & cuántas\\ \midrule
$\{\vac,X\}$ & indiscreta & 1\\
$\{\vac,\s{a},X\}$ & un punto abierto & 3\\
$\{\vac,\s{a,b},X\}$ & un par abierto & 3\\
$\{\vac,\s{a},\s{a,b},X\}$ & cadena & 6\\
$\{\vac,\s{a},\s{b,c},X\}$ & partición & 3\\
$\{\vac,\s{a},\s{b},\s{a,b},X\}$ & dos puntos abiertos y su unión & 3\\
$\{\vac,\s{a},\s{a,b},\s{a,c},X\}$ & un punto abierto y dos pares & 3\\
$\{\vac,\s{a},\s{b},\s{a,b},\s{a,c},X\}$ & dos puntos, dos pares & 6\\
$\Pt(X)$ & discreta & 1\\ \bottomrule
\end{tabular}\end{center}
$1+3+3+6+3+3+3+6+1=29$. Dos topologías de la misma fila se obtienen una de otra renombrando los puntos.

\section{El taller a la luz de la teoría}
\begin{center}\small\renewcommand{\arraystretch}{1.3}
\begin{longtable}{@{}p{2.3cm}p{4.3cm}p{8.8cm}@{}}\toprule
\textbf{Ejercicio} & \textbf{Qué evalúa} & \textbf{Clave y trampa}\\ \midrule\endhead
I y II, verificar & Axiomas de topología & Tablas de $\cup$ y $\cap$ de los pares no comparables. $\tau_2$ falla una
unión, $\tau_3$ una intersección. En II hay una errata en $\tau_1$.\\
I(a), II(a) & Cerrados & Complemento de cada abierto. Hay conjuntos abiertos y cerrados a la vez.\\
I(b), II(b) & Derivado & Usar el abierto mínimo y quitar el propio punto. Revisar \emph{todos} los puntos de $X$.\\
I(c), II(c) & Clausura & Menor cerrado que contiene al conjunto; comprobar con $A\cup A'$.\\
I(d), II(d) & Topología relativa & Intersecar cada abierto con $H$ y quitar repetidos.\\
I(e) & Densidad & $\cl{U}=X$. No confundir «denso» con «grande».\\
I(f), II(e) & Vecindades & Todos los superconjuntos de $U_x$. Convención de vecindad no necesariamente abierta.\\
I(g), II(f) & Interior & Unión de los abiertos contenidos. Puede ser $\vac$ o muy pequeño.\\
I(h), II(g) & Frontera ($\sigma$) & $\cl{A}\sm\inte{A}$. Notación por confirmar.\\
III.1--3 & Topologías pequeñas & Criterio de cadena o complementarios (4 términos); tres formas (5 términos).\\
III.4 & Comparación & Inclusión entre topologías; mismas cantidades de abiertos $\Rightarrow$ no comparables.\\
III.5 & Axiomas & Igual que la Parte I, con seis puntos.\\
III.6, III.7 & Topologías en conjuntos infinitos & Demostrar (T1)--(T3) con elementos; el abierto mínimo da las
vecindades. En III.7 falta $\vac$ en la definición.\\
III.8, III.9 & Topología generada & Primero intersecciones finitas, después uniones. En III.9 la sub-base no cubre $X$.\\
III.10 & Sub-bases & Familias que generan exactamente la topología dada.\\
III.11 & Sub-base en $\R$ & $[x-1,x]\cap[x,x+1]=\{x\}$: sale la discreta.\\
III.12 & Bases & Agrandar una base con abiertos sigue dando base. Errata en la conclusión.\\
III.13 & Axiomas & $A\subseteq B$, $B\subseteq A$ o $B=X\sm A$.\\
IV.1 & Continuidad global y puntual & Imágenes inversas; continua en $d$ pero no globalmente.\\
IV.2 & Función abierta & Contraejemplo $(-1,1)\mapsto[0,1)$.\\
IV.3 & Continua, abierta, cerrada & Tres pruebas independientes por función. $\tau_1$ no es topología; $X\not\homeo Y$.\\
IV.4 & Homeomorfismo & Contar abiertos para negar; construir la biyección para afirmar.\\ \bottomrule
\end{longtable}
\end{center}

% =====================================================================
\section{Recetario para el parcial}\label{sec:recetario}
\begin{center}\small\renewcommand{\arraystretch}{1.35}
\begin{longtable}{@{}p{4.3cm}p{11.1cm}@{}}\toprule
\textbf{Si piden\dots} & \textbf{Se hace\dots}\\ \midrule\endhead
¿Es topología? & (1) ¿Están $\vac$ y $X$? (2) Uniones e intersecciones de los pares no comparables. Un fallo basta
para decir que no.\\
Cerrados & Complemento de cada abierto.\\
Abierto mínimo $U_x$ & Intersección de los abiertos que contienen a $x$. Hacer la tabla al principio: sirve para todo
lo demás.\\
Vecindades $\N_x$ & Todos los conjuntos que contienen a $U_x$.\\
Derivado $A'$ & $x\in A'\iff(U_x\sm\{x\})\cap A\neq\vac$, para cada $x\in X$.\\
Clausura $\cl{A}$ & Menor cerrado que contiene a $A$; o $A\cup A'$; o $\{x:U_x\cap A\neq\vac\}$.\\
Interior $\inte{A}$ & Unión de los abiertos contenidos en $A$; o $\{x:U_x\subseteq A\}$.\\
Exterior & $\Int(X\sm A)=X\sm\cl{A}$.\\
Frontera & $\cl{A}\sm\inte{A}$. Comprobar $X=\inte{A}\sqcup\Fr(A)\sqcup\Ext(A)$.\\
¿Denso? & $\cl{D}=X$; o $D$ corta a todo abierto no vacío.\\
Topología relativa $\tau_H$ & $\{U\cap H:U\in\tau\}$, sin repetidos.\\
¿Comparables? & ¿Una lista de abiertos está contenida en la otra?\\
Topología generada & Intersecciones finitas (más $X$) $\to$ base; uniones (más $\vac$) $\to$ topología; verificar.\\
¿Base de $\tau$? & Miembros abiertos y todo abierto es unión de ellos.\\
¿Continua? & $f^{-1}(V)$ abierto en la salida para cada abierto $V$ de la llegada.\\
¿Continua en $p$? & Para cada abierto $V\ni f(p)$, ¿$f^{-1}(V)$ es vecindad de $p$? Finito: $f(U_p)\subseteq V_{f(p)}$.\\
¿Abierta? & $f(U)$ abierto en la llegada para cada abierto $U$ de la salida. Descarte: $f(X)$.\\
¿Cerrada? & $f(F)$ cerrado en la llegada para cada cerrado $F$ de la salida. Descarte: $f(X)$.\\
¿Homeomorfos? & Negar con un invariante (número de abiertos, de puntos abiertos\dots) o construir la biyección y
tabular $U\mapsto\varphi(U)$.\\
Demostrar que algo es topología (caso infinito) & (T1) mostrar que $\vac$ y $X$ cumplen la condición; (T2) tomar una
familia arbitraria $\{A_i\}$ y un elemento de la unión; (T3) tomar $A$, $B$ y un elemento de la intersección.\\ \bottomrule
\end{longtable}
\end{center}

% =====================================================================
\section{Errores frecuentes}
\begin{enumerate}[leftmargin=1.8em,itemsep=0.35em]
\item \textbf{Verificar solo uniones o solo intersecciones.} Hacen falta las dos (ejemplo: $\tau_3$ de la Parte~I).
\item \textbf{Pensar que cerrado es «no abierto».} Hay conjuntos con las dos propiedades y conjuntos sin ninguna.
\item \textbf{Olvidar el «distinto de $x$» en el derivado.} Un punto de $A$ puede no estar en $A'$.
\item \textbf{Buscar puntos de acumulación solo dentro de $A$.} Hay que revisar todo $X$.
\item \textbf{Confundir $A'$ con $\cl{A}$.} $\cl{A}=A\cup A'$; el derivado puede ser más pequeño o no contener a $A$.
\item \textbf{Creer que la frontera es «pequeña».} En espacios finitos puede ser casi todo $X$.
\item \textbf{Dar como $\N_x$ solo los abiertos} cuando la definición del curso admite vecindades no abiertas (o al
revés). Aclarar la convención en la respuesta.
\item \textbf{En la topología generada, unir antes de intersecar}, u olvidar agregar $X$ y $\vac$.
\item \textbf{Para la continuidad, calcular $f(U)$} en vez de $f^{-1}(V)$.
\item \textbf{Concluir «homeomorfos» porque hay una biyección} o porque coincide el número de abiertos. Hay que
exhibir la biyección y comprobar que lleva abiertos en abiertos en ambos sentidos.
\item \textbf{Decir «abierto» sin decir en qué espacio}, cuando hay subespacios de por medio.
\item \textbf{No leer el enunciado con cuidado.} Este taller tiene al menos cinco erratas; en un examen, si un dato
parece inconsistente, hay que señalarlo y explicar qué se asumió.
\end{enumerate}

% =====================================================================
\section{Autoevaluación (con respuestas)}\label{sec:auto}
Sea $X=\s{p,q,r,s}$ y $\tau=\{\vac,\s{p},\s{p,q},\s{p,r},\s{p,q,r},X\}$.
\begin{enumerate}[leftmargin=1.8em,itemsep=0.25em]
\item Compruebe que $\tau$ es una topología.
\item Halle los cerrados.
\item Halle el abierto mínimo de cada punto.
\item Para $A=\s{q,s}$ halle $A'$ y $\cl{A}$.
\item Halle $\cl{\s{p}}$ y $\s{p}'$. ¿Es $\s{p}$ denso?
\item Halle el interior de $\s{q,r,s}$ y el de $\s{p,q,s}$.
\item Para $B=\s{p,s}$ halle interior, clausura, frontera y exterior.
\item Halle $\N_s$ y $\N_q$.
\item Halle la topología relativa sobre $H=\s{q,r,s}$.
\item ¿Cuáles subconjuntos de $X$ son densos?
\item Sea $f:X\to X$ con $f(p)=p$, $f(q)=p$, $f(r)=q$, $f(s)=r$. ¿Es continua? ¿Abierta? ¿Cerrada?
\item Sea $\varphi:X\to X$ la función que intercambia $q$ con $r$ y fija $p$ y $s$. ¿Es un homeomorfismo?
\item Halle la topología $\tau^{\ast}$ sobre $X$ generada por $\C=\{\s{p,q},\s{q,r}\}$.
\item ¿Son homeomorfos $(X,\tau)$ y $(X,\tau^{\ast})$?
\item Verdadero o falso: (a) en un espacio discreto todo subconjunto es abierto y cerrado; (b) si $A$ es abierto,
entonces $A'=\vac$; (c) $\cl{A\cap B}=\cl{A}\cap\cl{B}$ siempre; (d) todo homeomorfismo es una función abierta; (e) la
unión de dos topologías es una topología; (f) toda biyección continua es homeomorfismo; (g) $\Fr(A)=\Fr(X\sm A)$
siempre; (h) el interior de un conjunto no vacío es no vacío.
\end{enumerate}

\begin{respuesta}
\textbf{\sffamily\color{verde!70!black}Respuestas.}
\begin{enumerate}[leftmargin=1.8em,itemsep=0.2em]
\item El único par no comparable es $\s{p,q}$, $\s{p,r}$: unión $\s{p,q,r}$, intersección $\s{p}$; ambos están.
\item $\cotau=\{\vac,\s{s},\s{q,s},\s{r,s},\s{q,r,s},X\}$.
\item $U_p=\s{p}$, $U_q=\s{p,q}$, $U_r=\s{p,r}$, $U_s=X$.
\item $A'=\s{s}$ (solo $U_s\sm\s{s}=\s{p,q,r}$ corta a $A$); $\cl{A}=\s{q,s}$, que es cerrado.
\item $\cl{\s{p}}=X$ y $\s{p}'=\s{q,r,s}$. Sí es denso.
\item $\Int\s{q,r,s}=\vac$; $\Int\s{p,q,s}=\s{p,q}$.
\item $\inte{B}=\s{p}$, $\cl{B}=X$, $\Fr(B)=\s{q,r,s}$, $\Ext(B)=\vac$.
\item $\N_s=\{X\}$; $\N_q=\{\s{p,q},\s{p,q,r},\s{p,q,s},X\}$.
\item $\tau_H=\{\vac,\s{q},\s{r},\s{q,r},H\}$.
\item Exactamente los que contienen a $p$ (8 conjuntos): todo abierto no vacío contiene a $p$.
\item No es continua: $f^{-1}(\s{p,r})=\s{p,q,s}\notin\tau$. Sí es abierta: las imágenes de los abiertos son
$\s{p},\s{p},\s{p,q},\s{p,q},\s{p,q,r}$. No es cerrada: $f(\s{s})=\s{r}$ no es cerrado.
\item Sí: es biyectiva y lleva $\tau$ sobre $\tau$ (intercambia $\s{p,q}$ con $\s{p,r}$ y fija los demás abiertos).
\item Base: $\{X,\s{q},\s{p,q},\s{q,r}\}$. $\tau^{\ast}=\{\vac,\s{q},\s{p,q},\s{q,r},\s{p,q,r},X\}$.
\item Sí: $\psi=(p\mapsto q,\ q\mapsto p,\ r\mapsto r,\ s\mapsto s)$ lleva $\s{p}\mapsto\s{q}$, $\s{p,q}\mapsto\s{p,q}$,
$\s{p,r}\mapsto\s{q,r}$, $\s{p,q,r}\mapsto\s{p,q,r}$, que es exactamente $\tau^{\ast}$.
\item (a) V. (b) F: en el ejercicio 5, $\s{p}$ es abierto y $\s{p}'\neq\vac$. (c) F: solo vale «$\subseteq$». (d) V.
(e) F. (f) F: la inversa puede no ser continua. (g) V. (h) F: ejercicio 6.
\end{enumerate}
\end{respuesta}

% =====================================================================
\section{Glosario de notación}
\begin{center}\small\renewcommand{\arraystretch}{1.3}
\begin{tabular}{@{}p{3.0cm}p{5.2cm}p{7.2cm}@{}}\toprule
\textbf{Aquí} & \textbf{Significado} & \textbf{Otras notaciones que aparecen en los textos}\\ \midrule
$\tau$ & topología (familia de abiertos) & $\mathcal{T}$, $\mathcal{O}$\\
$\cotau$ & familia de cerrados & $c_0\tau$ (taller), $\mathcal{F}$, $\mathcal{C}$\\
$\N_x$ & sistema de vecindades de $x$ & $N_x$ (taller), $\mathcal{V}(x)$, $\mathcal{U}_x$, $\mathcal{E}(x)$ (entornos)\\
$U_x$ & abierto mínimo de $x$ (espacios finitos) & $\mathrm{st}(x)$, $\uparrow\! x$\\
$A'$ & derivado & $A^{d}$, $\mathrm{Der}(A)$, $\mathrm{ac}(A)$\\
$\cl{A}$ & clausura o adherencia & $\mathrm{cl}(A)$, $\mathrm{Cl}\,A$, $A^{-}$, $\mathrm{adh}(A)$\\
$\inte{A}$, $\Int(A)$ & interior & $A^{\circ}$, $\mathrm{Int}\,A$\\
$\Ext(A)$ & exterior & $\mathrm{Ext}(A)$\\
$\Fr(A)$ & frontera & $\partial A$, $b(A)$, $\mathrm{Bd}\,A$; $\sigma(A)$ en el taller (por confirmar)\\
$\tau_H$ & topología relativa a $H$ & $\tau|_H$, $\tau/H$\\
$\B$, $\Sb$ & base, sub-base & $\mathfrak{B}$, $\mathcal{S}$\\
$\tau(\C)$ & topología generada por $\C$ & $\langle\C\rangle$\\
$X\homeo Y$ & $X$ homeomorfo a $Y$ & $X\approx Y$, $X\simeq Y$\\ \bottomrule
\end{tabular}\end{center}

% =====================================================================
\section*{Fuentes}
\addcontentsline{toc}{section}{Fuentes}
\markboth{Fuentes}{}
\begin{ojo}[Alcance de la verificación de fuentes]
El entorno en que se preparó esta guía permitía \emph{buscar} en la web, pero no abrir directamente las páginas ni los
PDF. Por eso: (1) los recursos en línea se consultaron a través de los resultados y resúmenes de búsqueda, no leyendo
el documento completo; (2) los libros impresos se citan por capítulo o sección según su índice, sin número de página,
y conviene confirmar la numeración en la edición que se tenga a mano; (3) todo resultado matemático de la guía se
puede comprobar de forma independiente: las demostraciones están escritas y los cálculos finitos se verificaron por
computadora. No se encontraron apuntes públicos del curso de la Universidad de Panamá, de modo que las convenciones de
notación propias del curso ($c_0\tau$, $\sigma$) quedan como supuestos.
\end{ojo}
\textbf{Libros de texto}
\begin{enumerate}[label={[\arabic*]},leftmargin=2.6em,itemsep=0.3em]\small\sloppy
\item S.~Lipschutz, \emph{Theory and Problems of General Topology}, Schaum's Outline Series, McGraw-Hill, Nueva York,
1965 (ed. en español: \emph{Topología general}). Capítulos «Topological spaces: definitions», «Bases and subbases» y
«Continuity and topological equivalence» (5, 6 y 7). Es, con mucha probabilidad, la fuente de los ejemplos de las
Partes I y IV del taller; la topología $\{X,\vac,\{a\},\{c,d\},\{a,c,d\},\{b,c,d,e\}\}$ es su ejemplo recurrente.
\item J.~R.~Munkres, \emph{Topology}, 2.ª ed., Prentice Hall, 2000. Capítulo 2, «Topological Spaces and Continuous
Functions»: §12 (espacios topológicos), §13 (bases), §16 (subespacios), §17 (cerrados y puntos límite), §18
(funciones continuas).
\item J.~L.~Kelley, \emph{General Topology}, Van Nostrand, 1955; reimpr. Springer, Graduate Texts in Mathematics 27.
Capítulo 1.
\item S.~Willard, \emph{General Topology}, Addison-Wesley, 1970; reimpr. Dover, 2004. Capítulos 2 y 3.
\item R.~Engelking, \emph{General Topology}, ed. revisada, Heldermann, Berlín, 1989. Capítulo 1.
\item G.~N.~Rubiano, \emph{Topología general (un primer curso)}, 2.ª ed., Universidad Nacional de Colombia, Bogotá,
2002.
\item M.~Macho Stadler, \emph{Topología}, apuntes del curso 2014--2015, Universidad del País Vasco (UPV/EHU).
\url{https://www.ehu.eus/~mtwmastm/Topologia1415.pdf}
\item F.~Cano, M.~Martín-Vega y B.~Molina-Samper, \emph{Topología: notas de un curso}, Universidad de Valladolid.
\url{https://uvadoc.uva.es/handle/10324/70838}
\item S.~A.~Morris, \emph{Topology Without Tears}, libro de acceso libre. \url{https://www.topologywithouttears.net}
\item L.~A.~Steen y J.~A.~Seebach, \emph{Counterexamples in Topology}, 2.ª ed., Springer, 1978. Ejemplos 5 (topología
de partición), 6 (par--impar) y 8 a 10 (punto particular), según los listados consultados.
\end{enumerate}
\textbf{Recursos en línea y artículos}
\begin{enumerate}[label={[\arabic*]},leftmargin=2.6em,itemsep=0.3em,start=11]\small\sloppy
\item OEIS Foundation, A000798: \emph{Number of different quasi-orders (or topologies, or transitive digraphs) with
$n$ labeled elements}. \url{https://oeis.org/A000798}
\item OEIS Foundation, A001930: \emph{Number of topologies, or transitive digraphs, with $n$ unlabeled nodes}.
\url{https://oeis.org/A001930}
\item Wikipedia, «Finite topological space». \url{https://en.wikipedia.org/wiki/Finite_topological_space}
\item Wikipedia, «Alexandrov topology». \url{https://en.wikipedia.org/wiki/Alexandrov_topology}
\item Wikipedia, «Kuratowski closure axioms». \url{https://en.wikipedia.org/wiki/Kuratowski_closure_axioms}
\item Wikipedia, «Subbase». \url{https://en.wikipedia.org/wiki/Subbase}
\item Wikipedia, «Comparison of topologies». \url{https://en.wikipedia.org/wiki/Comparison_of_topologies}
\item Wikipedia, «Particular point topology». \url{https://en.wikipedia.org/wiki/Particular_point_topology}
\item ProofWiki, «Definition: Odd-Even Topology». \url{https://proofwiki.org/wiki/Definition:Odd-Even_Topology}
\item B.~G.~Wachsmuth, \emph{Interactive Real Analysis}, §6.4, ejemplo de $x^2$ y la imagen de $(-1,1)$.
\url{https://mathcs.org/analysis/reals/cont/answers/imag1.html}
\item M.~Benoumhani, «The number of topologies on a finite set», \emph{Journal of Integer Sequences} 9 (2006),
artículo 06.2.6.
\item J.~Evans, \emph{Topology and Groups} (University College London), «2.07 Homeomorphisms».
\url{https://www.homepages.ucl.ac.uk/~ucahjde/tg/html/topsp07.html}
\item R.~Tinarrage, \emph{Topological Data Analysis} (EMAp/FGV), lección 2: homeomorfismos.
\url{https://raphaeltinarrage.github.io/files/EMAp/Lesson2.pdf}
\item Universidad de Extremadura, apuntes de \emph{Fundamentos Matemáticos}, tema 3 (notación $\mathrm{Fr}(A)$ y
$\mathrm{Ext}(A)$). \url{https://matematicas.unex.es/~trinidad/fmi/TTema3.pdf}
\item R.~P.~Stanley, «On the number of open sets of finite topologies», \emph{Journal of Combinatorial Theory} 10
(1971), 74--79.
\item N.~Bourbaki, \emph{Topologie générale}, capítulo 1 (estructuras topológicas).
\item K.~Ragnarsson y B.~E.~Tenner, «Obtainable sizes of topologies on finite sets», \emph{Journal of Combinatorial
Theory, Series A} 117 (2010), 138--151; arXiv:0802.2550.
\end{enumerate}

\end{document}
